\documentclass[margin,line]{res}
\usepackage[dvipsnames,table,xcdraw]{xcolor}
\usepackage{mdframed}
% \usepackage{newpxtext,newpxmath}
% \usepackage{txfonts}
\oddsidemargin -.5in
\evensidemargin -.5in
\textwidth=6.0in
\itemsep=0in
\parsep=0in
% if using pdflatex:
%\setlength{\pdfpagewidth}{\paperwidth}
%\setlength{\pdfpageheight}{\paperheight} 
\raggedbottom

% \pagestyle{empty}


\usepackage{xcolor}
\definecolor{linkblue}{HTML}{0000EE}
\usepackage[colorlinks=true, linkcolor=linkblue, urlcolor=linkblue, citecolor=linkblue]{hyperref}

\usepackage{bm}
\usepackage[symbol]{footmisc}

\newcommand\blfootnote[1]{%
  \begingroup
  \renewcommand\thefootnote{}\footnote{#1}%
  \addtocounter{footnote}{-1}%
  \endgroup
}
 
% Common Astronomy & Physics Journal Abbreviations
\newcommand{\aj}{AJ}                      % Astronomical Journal
\newcommand{\apj}{ApJ}                    % Astrophysical Journal
\newcommand{\apjl}{ApJL}                  % Astrophysical Journal Letters
\newcommand{\apjs}{ApJS}                  % Astrophysical Journal Supplement
\newcommand{\araa}{ARA\&A}                % Annual Review of Astronomy and Astrophysics
\newcommand{\aap}{A\&A}                   % Astronomy & Astrophysics
\newcommand{\aaps}{A\&AS}                 % Astronomy & Astrophysics Supplement
\newcommand{\aapr}{A\&ARv}                % Astronomy & Astrophysics Review
\newcommand{\mnras}{MNRAS}                % Monthly Notices of the RAS
\newcommand{\nat}{Nature}                 % Nature
\newcommand{\sci}{Science}                % Science
\newcommand{\pasp}{PASP}                  % Publications of the ASP
\newcommand{\pasj}{PASJ}                  % Publications of the Astronomical Society of Japan
\newcommand{\jcap}{JCAP}                  % Journal of Cosmology and Astroparticle Physics
\newcommand{\prd}{Phys.~Rev.~D}           % Physical Review D
\newcommand{\prl}{Phys.~Rev.~Lett.}       % Physical Review Letters
\newcommand{\physrep}{Phys.~Rep.}         % Physics Reports
\newcommand{\baas}{BAAS}                  % Bulletin of the AAS
\newcommand{\icarus}{Icarus}              % Icarus (planetary science)
\newcommand{\solphys}{Sol.~Phys.}         % Solar Physics
\newcommand{\ssr}{Space~Sci.~Rev.}        % Space Science Reviews
\newcommand{\apss}{Ap\&SS}                % Astrophysics and Space Science
\newcommand{\rmxaa}{RMxAA}                % Revista Mexicana de Astronomía y Astrofísica
\newcommand{\cjaa}{Chin.~J.~Astron.~Astrophys.} % Chinese Journal of Astronomy & Astrophysics
\newcommand{\zap}{ZAp}                    % Zeitschrift für Astrophysik (historic)
\newcommand{\na}{New~Astron.}             % New Astronomy
\newcommand{\nar}{New~Astron.~Rev.}       % New Astronomy Reviews
\newcommand{\actaa}{Acta Astron.}         % Acta Astronomica
\newcommand{\an}{Astron.~Nachr.}          % Astronomische Nachrichten
\newcommand{\aip}{AIP~Conf.~Proc.}        % AIP Conference Proceedings
\newcommand{\aspconf}{ASP~Conf.~Ser.}     % ASP Conference Series
\newcommand{\jqsrt}{J.~Quant.~Spec.~Radiat.~Transf.} % Journal of Quantitative Spectroscopy & Radiative Transfer
\newcommand{\memsai}{Mem.~Soc.~Astron.~Italiana} % Memorie della Società Astronomica Italiana
\newcommand{\memsaii}{Mem.~Soc.~Astron.~Italiana~Suppl.} % Supplements
\newcommand{\planss}{Planet.~Space~Sci.}  % Planetary and Space Science
\newcommand{\gca}{Geochim.~Cosmochim.~Acta} % Geochimica et Cosmochimica Acta
\newcommand{\natastro}{Nat.~Astron.}      % Nature Astronomy
\newcommand{\natastron}{Nat.~Astron.}     % Alternative alias
\newcommand{\natphys}{Nat.~Phys.}         % Nature Physics
\newcommand{\sovast}{Sov.~Astron.}        % Soviet Astronomy
\newcommand{\azh}{AZh}                    % Astronomicheskii Zhurnal
\newcommand{\ao}{Appl.~Opt.}              % Applied Optics
\newcommand{\jgr}{J.~Geophys.~Res.}       % Journal of Geophysical Research


\newlength{\cardwd}
\newlength{\cardgap}
\setlength{\cardgap}{4mm}      % horizontal gap between cards
 
% \setcardcols{N} : sets card width so N cards fill \linewidth exactly
\newcommand{\setcardcols}[1]{%
  \setlength{\cardwd}{%
    \dimexpr(\linewidth-\numexpr#1-1\relax\cardgap)/#1\relax}}
 
% \statcard{value}{label}
\newcommand{\statcard}[2]{%
  \begin{minipage}[t]{\cardwd}
    {\color{Maroon}\rule{\linewidth}{1.1pt}}\\[0.9mm]
    {\fontsize{19}{19}\selectfont\bfseries #1}\\[0.4mm]
    {\footnotesize\color{gray}\raggedright\textit{#2}\par}
  \end{minipage}\hfill}
 
\newcommand{\statcites}{28{,}000+}
\newcommand{\statrefd}{159}
\newcommand{\stath}{45}
\newcommand{\stati}{105}
\newcommand{\statfapapers}{10}


% \cardbreak : start a new row of cards, with vertical breathing room
\newcommand{\cardbreak}{\par\vspace{5.5mm}\noindent}
 
% \cardfill : pads out an incomplete final row so cards stay
%             left-aligned instead of being spread by \hfill
\newcommand{\cardfill}{\hspace*{\dimexpr\cardwd+\cardgap\relax}}

\newcommand{\mentee}[3]{%
  {\leftskip=5mm \noindent #1\hfill #2\par}%
  {\leftskip=11mm \noindent #3\par}%
  \vspace{1.5mm}}
\newcommand{\menteehead}[1]{\vspace{2mm}\noindent\textbf{#1}\par\vspace{0.5mm}}
 

\newlength{\awarddatewd}
\newlength{\awardtmp}
\newcommand{\award}[3][black]{%
  \settowidth{\awarddatewd}{\quad 2018, 2019}%
  \settowidth{\awardtmp}{\quad 2017--2021}%
  \ifdim\awardtmp>\awarddatewd \setlength{\awarddatewd}{\awardtmp}\fi
  \noindent
  \parbox[t]{\dimexpr\linewidth-\leftskip-\awarddatewd\relax}%
    {\raggedright\leavevmode\color{#1}#2\strut}%
  \hfill{\color{#1}#3}}
 

\renewcommand{\thefootnote}{\fnsymbol{footnote}}

\newenvironment{list1}{
  \begin{list}{\ding{113}}{%
      \setlength{\itemsep}{0in}
      \setlength{\parsep}{0in} \setlength{\parskip}{0in}
      \setlength{\topsep}{0in} \setlength{\partopsep}{0in} 
      \setlength{\leftmargin}{0.17in}}}{\end{list}}
\newenvironment{list2}{
  \begin{list}{$\bullet$}{%
      \setlength{\itemsep}{0in}
      \setlength{\parsep}{0in} \setlength{\parskip}{0in}
      \setlength{\topsep}{0in} \setlength{\partopsep}{0in} 
      \setlength{\leftmargin}{0.2in}}}{\end{list}}

\def\changemargin#1#2{\list{}{\rightmargin#2\leftmargin#1}\item[]}
\let\endchangemargin=\endlist 

\begin{document}

% \vspace{-.35cm}
\name{Joseph R. Farah \vspace*{.0in} \newline {\small\sc Miller Fellow of Astrophysics} \vspace*{.1in}}
\begin{resume}

\section{\sc Contact Information}
\vspace{.05in}
\begin{tabular}{@{}p{3in}p{4in}}
Miller Institute for Basic Research in Science & \textit{ORCID}: \href{https://orcid.org/0000-0003-4914-5625}{0000-0003-4914-5625} \\             
257 Campbell Hall & {\it E-mail:} \href{mailto:josephfarah@berkeley.edu}{\color{blue}josephfarah@berkeley.edu}\\
Berkeley, CA 94720  & {\it Website:} \href{http://jrfarah.com}{\color{blue}http://jrfarah.com} \\
USA & \textit{Citations}: \href{https://scholar.google.com/citations?hl=en&user=EVnW70sAAAAJ}{\color{blue}Google Scholar profile}\\
\end{tabular}

%\section{\sc Objective}
%Seeking a summer internship in {\bf computer vision} in either application or research

\section{\sc Research Interests}

I lead programs to measure the spins of compact objects across regimes, from infant millisecond magnetars in superluminous supernovae (SNe) to supermassive black holes resolved at the horizon scale. These measurements probe strong-field physics where no laboratory can reach and expose the engines driving the most energetic phenomena in the Universe.


{\sc Tags:}\par
\vspace{-0.4cm}
{\leftskip=3mm \noindent
\textit{Astrophysical Phenomena}: core-collapse SNe, superluminous SNe, black hole spacetimes, chirping SNe, photon rings \\
\textit{Observational Techniques}: very-long-baseline interferometry (VLBI), radio astronomy, optical photometry, optical spectroscopy \\
\textit{Computational \& Theoretical Methods}: general relativity, Bayesian modeling, machine/deep learning, geometric modeling, interferometric imaging, shock-cooling models, light curve modeling\par}


\vspace*{-1mm}
\section{\sc Education}
\label{sec:education}

{\bf University of California Santa Barbara}, Santa Barbara, CA, USA\hfill 2021--2026

\vspace{-.35cm}
{\em Ph.D. in Physics (completed)} \hfill 2026\\
{\em M.A. in Physics (completed) } \hfill 2023

{\bf University of Massachusetts Boston}, Boston, MA, USA\hfill 2017--2021

\vspace{-.35cm}
{\em B.S. in Physics, summa cum laude, with distinction}



% ============================================================
% PUBLICATION SUMMARY  
% ============================================================

% \vspace{-0.5cm}
\section{\sc {Publication} Summary}
\vspace*{-1.7mm}   % <-- TUNE THIS until the margin header lines up

\setcardcols{5}\noindent
\statcard{\statfapapers}{First-author}
\statcard{\statrefd}{Total}
\statcard{\statcites}{Citations}
\statcard{\stath}{$h$-index}
\statcard{\stati}{i10-index}

\begin{mdframed}[linecolor=Maroon]
\vspace{+.3cm}
\noindent {\bf \large Selected Publications}
\vspace{-.2cm}

{\leftskip=5mm
\noindent {\color{Maroon}\textbf{J. R. Farah}}, Prust, L. J., Howell, D. A., \textit{et al.}. 2026.
\textit{Lense-Thirring precessing magnetar engine drives a superluminous supernova}.
{\color{Maroon}\nat}, 651, \href{...}{321-325}. \\
{\color{gray} Reported the discovery of the first ``chirping'' supernova and its implications for the central engine hypothesis. Selected for the \href{https://www.nature.com/nature/volumes/651/issues/8105}{cover} of the March 12, 2026 issue of \textit{Nature}.}

\vspace*{+1mm}
\noindent {\color{Maroon}\textbf{J. R. Farah}}, Galison, P., Akiyama, K., \textit{et al.}. 2022.
\textit{Selective Dynamical Imaging of Interferometric Data}.
{\color{Maroon}\apjl}, 930, \href{https://iopscience.iop.org/article/10.3847/2041-8213/ac6615}{L18}. \\
{\color{gray} Seventh paper of the EHT's first Sagittarius~A$^\ast$ results series, with full Collaboration co-authorship. Method recognized by the LeRoy Apker Award, EHT Early Career Award.}
\par}
\vspace{0.35cm}
\end{mdframed}


% \hline



%%%%%%%%%%%%%%%%%%%%%

% \textit{GPA}: 3.9/4.0 
% \textit{GPA (physics)}: 3.909/4.0
\vspace{-.3cm}
\section{\sc Selected Work and Research Experience}


{\bf Miller Institute for Basic Research in Science}

\vspace{-.4cm}
{\em Miller Research Fellow, UC Berkeley --- Berkeley, CA} \hfill 2026--Present \\
\textit{Hosts: Prof.~Daniel Kasen, Prof.~Raffaella Margutti}
\end{changemargin}



{\bf Black Hole Initiative}

\vspace{-.4cm}
{\em Black Hole Explorer Fellow, Harvard University --- Cambridge, MA} \hfill 2026 \\
% 
% \begin{changemargin}{-0.0cm}{2.5cm} 
% I develop analytic and deep learning techniques for measuring physical properties from the black hole's photon ring. Along with the BHEX collaboration, I assist in the preparation and development of the BHEX spacecraft, an orbiting radio telescope which will interface with the EHT and use space VLBI to make high resolution images and movies of black holes.
% 
\textit{Supervisor: Dr.~Michael Johnson}
\end{changemargin}




{\bf Las Cumbres Observatory}\hfill 2021--2026


\vspace{-.4cm}
{\em NSF Graduate Research Fellowship --- Goleta, CA} \\
%
% \begin{changemargin}{-0.0cm}{2.5cm} 
% I use the LCO global network of telescopes to perform optical and spectroscopic follow-up of a wide variety of supernovae and transients. Using rapid follow-up observations combined with new shock cooling physics, I seek to better characterize and understand the mass stripping mechanism of Type IIb supernovae, and the underlying mechanisms powering superluminous supernovae. Additionally, I help with infrastructure maintenance and development for the robotic LCO platform.
%
\textit{Supervisor: Prof.~Andy Howell}
\end{changemargin}



{\bf Harvard-Smithsonian Center for Astrophysics} \hfill 2018--2021

\vspace{-.4cm}
{\em Smithsonian Fellowship with the Event Horizon Telescope --- Cambridge, MA}\\
%
% \begin{changemargin}{-0.0cm}{2.5cm} 
% I image supermassive black holes and assist in observations using the Event Horizon Telescope in order to help take the first ever pictures of black holes. I explored alternative parametric representations of shadows in the Kerr metric and non-standard GR theories, and methods constraining the angular momentum of rotating black holes in EHT data using image- and Fourier-domain feature extraction methods. As an undergraduate, I was the first and only junior collaborator of the EHT Consortium. 
%
\textit{Supervisor: Dr.~Michael Johnson}
\end{changemargin}

{\bf Undergraduate Research, UMass Boston and Harvard University}\hfill 2017--2021

\vspace{-.4cm}
{\em Quantum computing with Dr.~Alioscia Hamma (2019--2021); QBism as an Oracle Fellow with Dr.~Christopher Fuchs (2017--2019); ATLAS Micromegas detector instrumentation with Dr.~Melissa Franklin, Harvard Laboratory for Particle Physics and Cosmology (2017--2018)}

% {\bf Harvard School for Engineering and Applied Sciences}\hfill Summer 2016

% \vspace{-.4cm}
% {\em Summer Research Intern mentored by Dr.~Chinwendu Enyioha --- Cambridge, MA}

% {\bf Tufts Department of Astrophysics}\hfill Summer 2015

% \vspace{-.4cm}
% {\em Summer Research Intern mentored by Dr.~Anna Sajina --- Medford, MA}



%%%%%%%%%%%%%%
\vspace{+.2cm}
\section{\sc Leadership and \\ Collaboration Roles}
 
{\bf Black Hole Explorer (BHEX)}
 
\vspace{-.4cm}
{\em Deputy Science Lead, Goal I (Photon Rings)} \hfill 2026--Present\\ 
{\em Core Member} \hfill 2020--Present\\
{\color{gray} Contributing to the science case of the proposed NASA space-VLBI mission since 2020. Lead development of its black hole spin-inference capabilities (i.e., photon ring analyses, cross-validation, and simulated data challenges) with extended line of first-author BHEX papers. Advise the PI on development of the BHEX Student Collaboration program.}
 
\vspace*{-2mm}
{\bf CHIRPS Collaboration}\hfill 2026--Present
 
\vspace{-.4cm}
{\em Founder and Coordinator}\\
{\color{gray} \textit{CHaracterization and Inference of Rotating, magnetar-Powered SLSNe.} Multi-institution collaboration (17 members across 9 institutions) for the discovery and modeling of chirping superluminous supernovae. Operates a public real-time chirp-detection filter on the ANTARES broker for Rubin/LSST alerts and leads the development of theoretical interpretations. {\color{blue}\href{https://chirps.space}{Website.}}} 
 
\vspace*{-2mm}
{\bf Research Group, UC Berkeley}\hfill 2026--Present
 
\vspace{-.4cm}
{\em Group Lead}\\
{\color{gray} Supervise 5 student researchers at UC Berkeley, UC Santa Barbara, and Northwestern/CIERA on BHEX spin inference and the CHIRPS program.} %% CHECK N
 
\vspace*{-2mm}
{\bf UltraViolet EXplorer (UVEX)}\hfill 2026--Present
 
\vspace{-.4cm}
{\em Collaboration Member, Supernova Science}\\
{\color{gray} Assisting with development of the supernova observing strategy, including the cadence required to constrain early temperature evolution.}
 


 
\vspace*{-2mm}
{\bf Global Supernova Project}\hfill 2021--Present
 
\vspace{-.4cm}
{\em Builder}\\
{\color{gray} Contributed to core collaboration infrastructure (including the Supernova Exchange) and performed data reduction for collaboration papers: 1000+ optical spectra and dozens of multi-band light curves. Can trigger Las Cumbres 1-m and 2-m photometric and spectroscopic follow-up.}

\vspace*{-2mm}
{\bf Event Horizon Telescope Collaboration}\hfill 2018--Present 
 
 \vspace{-.4cm}
{\em Collaboration Member}\\
{\color{gray} Joined as the collaboration's first undergraduate member. Contributed to imaging, modeling, and theory for the first images of M87$^\ast$ (2019) and Sgr~A$^\ast$ (2022). Led the selective dynamical imaging paper in the 2022 Sgr~A$^\ast$ results (\apjl, 930, L18), recognized with the EHT Early Career Award.}

\newpage
\vspace{+.2cm}

\section{\sc {\color{Maroon} Selected\\ Honors} and \\ Awards}
\vspace*{0.67cm}
 
\award{{\color{Maroon}Miller Fellowship, Miller Institute for Basic Research in Science, Berkeley}}{2026--2029} \\
{\color{gray}\textsc{Award amount}: \$354,000}
 
\vspace*{-2mm}
\award{Bhaumik Prize Postdoctoral Fellowship, UCLA (declined)}{2026} \\
{\color{gray}\textsc{Award amount}: \$276,000}
 
\vspace*{-2mm}
\award{UC Santa Barbara Mananya Tantiwiwat Award}{2025}
% {\color{gray}\textsc{Award amount}: \$1500}
 
\vspace*{-2mm}
\award{Broida-Hirschfelder Fellowship}{2025} \\
{\color{gray}\textsc{citation}: \textit{In recognition of his project uncovering the secrets of the strangest object in the universe, the black hole photon ring.}}
% {\color{gray}\textsc{Award amount}: \$10,834}
 
\vspace*{-2mm}
\award{$1^\textrm{st}$ Place Award, UC Santa Barbara Grad Slam}{2025} \\
{\color{gray}\textit{Round 1 (Physics and Chemistry)}}
 
\vspace*{-2mm}
\award{UC Santa Barbara GSA Excellence in Teaching Award (nominated)}{2025}
 
\vspace*{-2mm}
\award{UC Santa Barbara Academic Senate Outstanding Teaching Assistant Award (nominated)}{2024}
 
\vspace*{-2mm}
\award{Event Horizon Telescope Early Career Award}{2022} \\
{\color{gray}\href{https://twitter.com/ehtelescope/status/1602826998346551296}{\textsc{citation}}: \textit{For his leadership and contribution to the dynamic imaging of the interferometric data used to image Sgr A*.}}
 
\vspace*{-2mm}
\award[Maroon]{LeRoy Apker Award}{2021} \\
{\color{gray}\href{https://www.aps.org/funding-recognition/award/leroy-apker-award}{\textsc{citation}}: \textit{For the invention of the selective dynamical imaging method, with applications for studying rapidly-varying black holes.}}
% {\color{gray}\textsc{Award amount}: \$10,000}
 
\vspace*{-2mm}
\award{UC Santa Barbara James and Mary Jo Hartle Graduate Fellowship}{2021}
% {\color{gray}\textsc{Award amount}: \$2,463}
 
\vspace*{-2mm}
\award[Maroon]{National Science Foundation Graduate Research Fellowship}{2021--2026} \\
{\color{gray}\textsc{Award amount}: \$138,000}
 
\vspace*{-2mm}
\award[Maroon]{Breakthrough Prize in Fundamental Physics (co-recipient)}{2020}
 
\vspace*{-2mm}
\award[Maroon]{Barry M. Goldwater Scholarship}{2019}
% {\color{gray}\textsc{Award amount}: \$15,000}
 
\vspace*{-2mm}
\award{National Science Foundation Diamond Achievement Award$^\ddag$}{2019}
 
\vspace*{-2mm}
\award{Smithsonian Fellowship}{2018, 2019}
 
\vspace*{-2mm}
\award{Oracle Fellowship}{2018, 2019}
 
\vspace*{-2mm}
\award{AJAS/AAAS Lifetime Fellowship}{2018}
 
\vspace*{-2mm}
\award{\textbf{Selected undergraduate scholarships}: Chancellor's Scholarship (2017--2021), Alton J.~Brann Endowed Scholarship (2019), Northrop Grumman Scholarship (2020), Arthur W.~Martin III Scholarship (2021)}{2017--2021}
 

% \vspace*{-2mm}
% $2^\textrm{nd}$ Place Award, Massachusetts State Science and Engineering Fair \hfill 2017

% \vspace*{-2mm}
% $1^\textrm{st}$ Place Award, Massachusetts State Science and Engineering Fair \hfill 2016


% \section{\sc Academic\\ Service}
\vspace{+.0cm}
\section{\sc Invited Talks \\ and Colloquia}

% (last updated: August 2026)

A new theoretical framework for chirping supernovae.  Montana State University, Bozeman, MT. October 2026.

\vspace*{-2mm}
Precessing magnetar central engine drives a superluminous supernova. Galactic Center Group, UC Los Angeles, Los Angeles, CA. April 2026.

\vspace*{-2mm}
The first 4 years of SN 1993J revisited. UC Berkeley, Berkeley, CA. September 2025.

\vspace*{-2mm}
Lense-Thirring precessing magnetar engine drives a superluminous supernova. Multi-RAPTOR colloquium, UC Berkeley, Berkeley, CA. September 2025.

\vspace*{-2mm}
Inferring black hole spin from images of the photon ring. VandyGRAF Initiative, Vanderbilt University, Nashville, TN. September 2024.

\vspace*{-2mm}
What Lies Within: Imaging the Galactic Center. April APS Meeting, New York, NY. Invited session speaker. April 2022.

\vspace*{-2mm}
The First Image of a Black Hole. Skyscrapers, Inc., Seagrave Memorial Observatory, North Scituate, RI. 2021.



\vspace{+.2cm}
\section{\sc Seminars and Contributed Talks}

A new theoretical framework for chirping supernovae. Multi-RAPTOR seminar, UC Berkeley, Berkeley, CA. October 2026.

\vspace*{-2mm}
Spin inference of compact objects across regimes. Black Hole Initiative, Cambridge, MA. August 2026.

\vspace*{-2mm}
Precessing magnetar central engine drives a superluminous supernova. METAL group, University of Virginia, Charlottesville, VA. August 2026.

\vspace*{-2mm}
Spin inference of compact objects across regimes. Greene research group, Princeton University, Princeton, NJ. July 2026.

\vspace*{-2mm}
Precessing magnetar central engine drives a superluminous supernova. Narayan group, Center for Astrophysics $|$ Harvard \& Smithsonian, Cambridge, MA. July 2026.

% \vspace*{-2mm}
% Presentation presentation: a seminar on seminars. Workshop, Las Cumbres Observatory, Goleta, CA. January 2026.

\vspace*{-2mm}
General relativity meets a supernova for the first time. Las Cumbres Observatory colloquium, Goleta, CA. January 2026. {\color{blue}\href{https://lco.global/seminar/jfarah/}{Event URL.}}

\vspace*{-2mm}
Lense-Thirring precessing magnetar engine drives a superluminous supernova. KIPAC, Stanford University, Stanford, CA. November 2025.

\vspace*{-2mm}
Machine- and deep-learning-driven angular momentum inference from Black Hole Explorer observations of the $n=1$ photon subring. \textit{Topics in Gravitational Physics}, APS Global Physics Summit. March 2025. {\color{blue}\href{https://summit.aps.org/events/APR-C20/8}{Event URL.}}

\vspace*{-2mm}
Current and future impacts of survey capabilities on the stripped-envelope problem. Kavli Institute for Theoretical Physics, Santa Barbara, CA. February 2025.

\vspace*{-2mm}
Shock-cooling constraints via early-time observations of the Type IIb SN 2022hnt. Kavli Institute for Theoretical Physics, Santa Barbara, CA. February 2025.

\vspace*{-2mm}
Selective dynamical imaging of interferometric data and the second image of a black hole. \textit{Astro Tea}, University of California, Santa Barbara, CA. June 2022.

\vspace*{-2mm}
Seeing the unseeable: the first image of a black hole. University of Massachusetts Boston, Boston, MA. May 2019.

\vspace*{-2mm}
Imaging a black hole with the Event Horizon Telescope. Harvard University, Cambridge, MA. April 2019.

% \vspace{+.2cm}
% \section{\sc Internal and Institutional Talks}





\vspace{+.2cm}
\section{\sc Teaching and Mentoring}

Invited workshop creator and lecturer, \textit{Presentation presentation: a seminar on seminars}, Las Cumbres Observatory, Goleta, CA. January 2026.

\vspace*{-2mm}
Death Star Seminar, founder and lecturer. Informal graduate seminar series on supernova theory, Las Cumbres Observatory, Goleta, CA. Designed curriculum and delivered lectures. 2026.

\vspace*{-2mm}
Creator and lecturer, online derivation series, \textit{Astraveo}. Blackboard lectures working through core astrophysics and general science results from first principles. 2025--Present. 

\vspace*{-2mm}
Teaching Assistant, PHYS 23 (Electrostatics and waves), UC Santa Barbara. Fall 2024.\par
\vspace{-0.4cm}
{\leftskip=4mm \noindent
{UC Santa Barbara GSA Excellence in Teaching Award (nominated)} \\
{UC Santa Barbara Academic Senate Outstanding Teaching Assistant Award (nominated)}\par}

\vspace*{-2mm}
Guest Lecturer. Astronomy 121, University of Massachusetts Boston, Boston, MA. 2019 (x2), 2020 (x2).

{\setlength{\parskip}{0pt}
 
\vspace*{0.3cm}
\menteehead{Graduate research mentoring}
 
\mentee{Nicole Flors (CIERA/Northwestern)}{2026--Present}
  {Project: Simulation of magnetar+disk system with HAM-R}
 
\menteehead{Undergraduate research mentoring}
 
\mentee{Eyouel Abate (UC Berkeley)}{2026--Present}
  {Project: Development of automated CHIRPS filter on Rubin infrastructure}
 
\mentee{Austin Mei (UC Berkeley)}{2026--Present}
  {Project: Frequency-domain methods for identification of periodicity in supernova light curves}
 
\mentee{Mai Haller (UC Berkeley)}{2026--Present}
  {Project: Investigation of contrast ratios of SLSNe to surface brightness of hosts}
 
\mentee{Courtney Duong (UC Santa Barbara)}{2024--Present}
  {Project: Convolutional neural network for extraction of the $n=1$ black hole photon ring (Published in \textit{ApJ}) \\
   Award: \textit{Gene and Susan Lucas Undergraduate Research Scholarship} \\
   Award: \textit{Edison STEM Research Scholarship}}
 
\mentee{Franklin Myhre (UC Santa Barbara)}{2024--Present}
  {Project: Development of a \texttt{python} library for automated image-domain subring extraction}
 
\mentee{Sanjit Masanam (UC Santa Barbara)}{2024--2026}
  {Award: \textit{2025 Goldwater Scholarship (nominated)} \\
   Award: \textit{2025 Ernest F. Hollings Scholarship}}
 
\mentee{Songgun Lee (UC Santa Barbara)}{2024--2026}
  {\textit{Co-founder at Masterminding}}
 
\mentee{Tazzy Imbabi (UC Santa Barbara)}{2024}
  {Project: Exploration of selective dynamical imaging applications to BHEX}
 
}


% \vspace{+.1cm}
\section{\sc Service}

{\bf Referee:}\par
\vspace{-0.4cm}
{\leftskip=3mm \noindent
\textit{The Astrophysical Journal} \\
\textit{Astronomy \& Astrophysics}\par}

\vspace*{-2mm}
Internal referee, Event Horizon Telescope Collaboration.

\vspace*{-2mm}
Invited presenter, Las Cumbres Observatory Science Advisory Committee, Goleta, CA. Invited to present on the observatory's high-energy astrophysics science program. August 2025.

\vspace*{-2mm}
Local Organizing Committee, EHT Collaboration Meeting. 2020.

\vspace*{-2mm}
Local Organizing Committee, 2nd EHT Imaging Workshop, Black Hole Initiative, Harvard University, Cambridge, MA. July 2018.

\vspace*{-2mm}
Founder and organizer, \textit{Research@UMB}, University of Massachusetts Boston. Connected undergraduates seeking research experience with faculty through faculty lightning-talk sessions and direct student--faculty matching. 2019. %% CHECK "Founder" vs "Organizer", year


\vspace{+.2cm}
\section{\sc Outreach and Public Engagement}

% \vspace*{-2mm}
Co-founder and host, \textit{Astraveo}, science-communication YouTube channel (with D.~James). 2023--Present. {\color{blue}\href{https://www.youtube.com/@astraveo}{Channel URL.}}\par %% CHECK year, handle
\vspace{-0.4cm}
{\leftskip=4mm \noindent
{\color{gray} 14,000+ subscribers; most recent video (including an on-camera blackboard lecture) 150,000+ views; 4+ videos with 25,000+ views each, 100+ videos total; research vlog, documentaries.}\par} 


General relativity beats the heart of a dying star. Public lecture, Santa Barbara Museum of Natural History. March 2026. {\color{blue}\href{https://www.sbnature.org/visit/calendar/9132/general-relativity-beats-the-heart-of-a-dying-star}{Event URL.}}

\vspace*{-2mm}
Extragalactic science with the Vera C. Rubin Observatory. Public lecture, Santa Barbara Museum of Natural History. August 2025. {\color{blue}\href{https://www.sbnature.org/visit/calendar/8760/new-frontiers-in-astronomy-powered-by-the-rubin-observatory}{Event URL.}}

\vspace*{-2mm}
The Black Hole Explorer and the edge of the universe. Public lecture, Santa Barbara Museum of Natural History. April 2025. {\color{blue}\href{https://www.sbnature.org/visit/calendar/8642/the-black-hole-explorer-and-the-edge-of-the-universe}{Event URL.}}

\vspace*{-2mm}
The Black Hole Explorer. \textit{Astronomy on Tap}, Santa Barbara, CA. December 2024. {\color{blue}\href{https://lco.global/aot/december-5-2024/}{Event URL.}}

\vspace*{-2mm}
What does an astronomer do? Canalino Elementary School, Carpinteria, CA. October 2024.


\vspace*{-2mm}
The biggest telescopes! Aliso Elementary School, Carpinteria, CA. October 2023.


\vspace*{-2mm}
Imaging the black hole at our galaxy's center. Public lecture, Santa Barbara Museum of Natural History. May 2023. {\color{blue}\href{https://www.sbnature.org/visit/calendar/7329/free-public-astronomy-talk-imaging-the-black-hole-at-our-galaxys-center}{Event URL.}}

 

\vspace*{-2mm}
LCO's top 9 tips for taking your best black hole photo. \textit{Astronomy on Tap}, Santa Barbara, CA. November 2022. {\color{blue}\href{https://astronomyontap.org/event/astronomy-on-tap-santa-barbara-2/}{Event URL.}}

\vspace*{-2mm}
AskScience AMA Series: Event Horizon Telescope results on the Galactic Center. \textit{Reddit}. May 2022. {\color{blue}\href{https://old.reddit.com/r/askscience/comments/unyg77/askscience_ama_series_were_event_horizon/}{Article URL.}}

\vspace*{-2mm}
LCO's top 9 tips for taking your best black hole photo. Honors and AP physics classes, Medford High School, Medford, MA. 2022.

\vspace{+.2cm}
% \section{\sc Filmography}

% Galison, P. L. (Director). 2021. \textit{Black Holes: The Edge of All We Know.} Sandbox Films. {\color{blue}\href{https://www.netflix.com/title/81343342}{View on Netflix.}}

% \vspace*{-2mm}
% Galison, P. L. (Director). 2019. \textit{Portrait of a Shadow.} Sandbox Films.

\vspace{+.2cm}
\section{\sc Selected Media \\ and Press}
 
{\bf Appearances}\par
\vspace{-0.4cm}
{\leftskip=3mm \noindent
Guest, \textit{Nature Podcast}, on the chirping supernova SN~2024afav. March 2026. {\color{blue}\href{https://youtu.be/ZfJA6FJICns?t=691}{Listen.}} \\ %% CHECK episode title and date
Participant, Event Horizon Telescope press conference unveiling the first image of Sgr~A$^\ast$, Washington, D.C. May 2022. \\ %% CHECK your role (panelist/speaker?)
Galison, P.~L. (Director). 2021. \textit{Black Holes: The Edge of All We Know.} Sandbox Films. {\color{blue}\href{https://www.netflix.com/title/81343342}{View on Netflix.}} \\
Galison, P.~L. (Director). 2019. \textit{Portrait of a Shadow.} Sandbox Films.\par}
 
\vspace*{-1mm}
{\bf Coverage}\par
\vspace{-0.4cm}
{\leftskip=3mm \noindent
\textit{SN 2024afav (Nature, 2026):} %
\href{https://www.reuters.com/science/astrophysicists-figure-out-what-caused-super-bright-supernova-2026-03-11/}{Reuters},
\href{https://www.thetimes.com/uk/science/article/astronomers-magnetar-birth-first-2j8drllkw}{\textit{The Times}},
\href{https://www.independent.co.uk/space/supernova-superluminous-cause-neutron-star-b2936619.html}{\textit{The Independent}},
\href{https://www.science.org/content/article/universe-s-brightest-stellar-explosions-may-be-powered-highly-magnetic-neutron-stars}{\textit{Science}},
\href{https://www.scientificamerican.com/article/the-universes-brightest-supernovas-are-turbocharged-by-newborn-magnetars/}{\textit{Scientific American}},
\href{https://www.nature.com/articles/d41586-026-00490-3}{\textit{Nature} News \& Views},
\href{https://www.youtube.com/watch?v=mSoIMgFTyPc}{\textit{Nature} Video},
\href{https://www.sciencenews.org/article/superluminous-supernova-magnetar}{\textit{Science News}},
\href{https://arstechnica.com/science/2026/03/magnetars-drag-spacetime-to-power-superluminous-supernovae/}{\textit{Ars Technica}},
\href{https://www.smithsonianmag.com/smart-news/astronomers-witnessed-the-birth-of-a-magnetar-for-the-first-time-it-explains-the-mysterious-flickering-of-an-ultra-bright-supernova-180988343/}{\textit{Smithsonian}},
\href{https://www.popsci.com/science/first-magnetar-birth/}{\textit{Popular Science}},
\href{https://skyandtelescope.org/astronomy-news/super-bright-supernovae-are-magnetar-birth-cries/}{\textit{Sky \& Telescope}},
\href{https://www.skyatnightmagazine.com/news/sn-2024afav-birth-of-magnetar}{\textit{BBC Sky at Night}},
\href{https://www.space.com/astronomy/stars/astronomers-witness-colossal-supernova-explosion-create-one-of-the-most-magnetic-stars-in-the-universe-for-the-first-time}{\textit{Space.com}},
\href{https://livescience.com/space/astronomy/scientists-witness-birth-of-one-of-the-universes-strongest-magnets-for-the-first-time-thanks-to-a-general-relativity-magic-trick}{\textit{Live Science}},
\href{https://www.sciencealert.com/strange-chirp-may-reveal-what-powers-the-brightest-supernovae-in-the-universe}{\textit{ScienceAlert}},
\href{https://www.iflscience.com/first-direct-evidence-of-magnetar-formation-from-a-superluminous-supernova-82832}{\textit{IFLScience}};
international coverage including a print feature in \textit{T\"ahdet ja avaruus} (Finland), %% CHECK issue/date
\href{https://www.spektrum.de/news/die-hellsten-supernovae-des-universums-entstehen-durch-magnetare/2314478}{\textit{Spektrum der Wissenschaft}} (Germany),
\href{https://www.agenciasinc.es/Noticias/Una-supernova-emite-extranas-ondulaciones-de-luz-debido-a-un-magnetar-en-su-interior}{Agencia SINC} (Spain),
\href{https://www.media.inaf.it/2026/03/16/lense-thirring-magnetar-sn-2024afav/}{\textit{Media INAF}} (Italy),
\href{https://www.derstandard.at/story/3000000312039/extreme-raumzeitverzerrung-erstmals-in-supernova-nachgewiesen}{\textit{Der Standard}} (Austria), and
\href{https://www.lanacion.com.ar/ciencia/algo-raro-esta-pasando-cientificos-captan-un-extrano-chirrido-que-proviene-del-espacio-profundo-nid03052026/}{\textit{La Naci\'on}} (Argentina);
YouTube features by \href{https://www.youtube.com/watch?v=tlmj5ThgWbY}{Astrum} and \href{https://youtu.be/eT5INRJYzwA?si=0Seh0etgvrDRXmDl}{Anton Petrov}. 
Press releases by
\href{https://news.berkeley.edu/2026/03/11/astronomers-capture-birth-of-a-magnetar-confirming-link-to-some-of-universes-brightest-exploding-stars/}{UC Berkeley},
\href{https://news.ucsb.edu/2026/022435/ucsb-researcher-bridges-worlds-general-relativity-and-supernova-astrophysics}{UC Santa Barbara}, and
\href{https://lco.global/news/mysterious-power-source-for-superluminous-supernovae-revealed-by-general-relativity/}{Las Cumbres Observatory}. \\[1mm]
\textit{Selective dynamical imaging of Sgr~A$^\ast$ (2022):} %
profiles in the \href{https://www.independent.com/2022/05/19/santa-barbara-astronomer-bags-a-black-hole/}{\textit{Santa Barbara Independent}} and on
\href{https://www.kclu.org/local-news/2022-05-12/south-coast-scientist-played-a-critical-role-in-capturing-the-first-image-of-our-galaxys-black-hole}{KCLU radio};
\href{https://www.eurekalert.org/news-releases/952599}{UC Santa Barbara/LCO press release}. \\[1mm]
\textit{Event Horizon Telescope, undergraduate work (2019--2021):} %
profile in \href{https://www.symmetrymagazine.org/article/drag-racing-and-black-hole-physics}{\textit{Symmetry} magazine} (2020);
\href{https://pweb.gws.cfa.harvard.edu/news/rising-eht-scientist-awarded-prestigious-physics-award}{Center for Astrophysics $|$ Harvard \& Smithsonian} and
\href{https://spmt.umb.edu/news/2021/umass-boston-alum-wins-prestigious-national-academ}{UMass Boston} coverage of the Apker Award (2021).\par}

\phantomsection\label{sec:firstauthor}


\vspace{+.8cm}
\section{{\sc Publications (first-author)  \\ \vspace{0.3cm} {\footnotesize \textbf{\color{Maroon} Stats (all papers)}: \\Citations: \statcites \\ \vspace{0.1cm} \noindent Ref'd Papers: \statrefd \vspace{0.1cm} \\ $h$-index: \stath \\ i10-index: \stati}} \\ }


10. \textbf{J. R. Farah}, Prust, L. J., Yu, H. 2026. \textit{A universal relation predicts the observed reversal of a supernova chirp}. arXiv:\href{https://arxiv.org/abs/2610.39240}{2610.39240}.

% \textbf{NB:} Please see full list of publications on page 6
9. \textbf{J. R. Farah}, Palumbo, D. C. M., Johnson, M. D., Chang, D. O. 2026. \textit{Spin inference with the Black Hole Explorer. I. Fisher information matrix forecast}. arXiv:\href{https://arxiv.org/abs/2608.23672v1}{2608.23672}.

8. \textbf{J. R. Farah}, Lupsasca, A., Quataert, E., Johnson, M. D. 2026. \textit{Interferometric inference of black hole spin from photon ring size and brightness}. ApJ, 1003, \href{https://iopscience.iop.org/article/10.3847/1538-4357/ae6116}{63}. 
%arXiv:\href{https://arxiv.org/abs/2509.23628}{2509.23628}.

7. \textbf{J. R. Farah}, Prust, L. J., Howell, D. A., \textit{et al.}. 2026. \textit{Lense-Thirring precessing magnetar engine drives a superluminous supernova}. Nature, 651, \href{https://www.nature.com/nature/volumes/651/issues/8105}{321-325}. \textbf{NB:} selected for \href{https://www.nature.com/nature/volumes/651/issues/8105}{cover} of the March 12, 2026 issue. 
%\href{https://arxiv.org/abs/2509.08051}{arXiv:2509.08051}.

6. \textbf{J. R. Farah}, Howell, D. A., Hiramatsu, D., \textit{et al.}. 2026. \textit{When IIb Ceases To Be: Bridging the Gap Between IIb and Short-plateau Supernovae}. ApJ, 998, \href{https://iopscience.iop.org/article/10.3847/1538-4357/ae3a71}{321}. 
%\href{https://arxiv.org/abs/2509.12470}{arXiv:2509.12470}.

5. \textbf{J. R. Farah}, Prust, L. J., Terreran, G., \textit{et al.}. 2025. \textit{The First 4 Years of SN 1993J Revisited: Geometric $m$-ring Modeling of the Radio Shell with Closure Quantities Only}. Submitted to ApJ, positive referee report received. arXiv e-prints, \href{https://arxiv.org/abs/2509.20601}{arXiv:2509.20601}.

4. \textbf{J. R. Farah}, Davelaar, J., Palumbo, D., Johnson, M., Delgado, J. 2025. \textit{Machine and Deep Learning-driven Angular Momentum Inference from BHEX Observations of the n = 1 Photon Ring}. \apj, 983, \href{https://iopscience.iop.org/article/10.3847/1538-4357/adbbe3}{185}. 
%\href{https://arxiv.org/abs/2411.01060}{arXiv:2411.01060}.

3. \textbf{J. R. Farah}, Howell, D. A., Terreran, G., \textit{et al.}. 2025. \textit{Shock-cooling Constraints via Early-time Observations of the Type IIb SN 2022hnt}. \apj, 984, \href{https://iopscience.iop.org/article/10.3847/1538-4357/adad5a}{60}. 
%\href{https://arxiv.org/abs/2501.17221}{arXiv:2501.17221}.


2. \footnote[$^S$]{\noindent $^S$This was the seventh paper of the 10-paper first Galactic Center results series from the Event Horizon Telescope, presenting the second image of a black hole and the first horizon-scale images of Sgr A$^\ast$.}\textbf{J. R. Farah}, Galison, P., Akiyama, K., \textit{et al.}. 2022. \textit{Selective Dynamical Imaging of Interferometric Data}. \apjl, 930, \href{https://iopscience.iop.org/article/10.3847/2041-8213/ac6615}{L18}. \href{https://arxiv.org/abs/2409.08321}{arXiv:2409.08321}.

1. \textbf{J. R. Farah}, Pesce, D. W., Johnson, M. D., Blackburn, L. 2020. \textit{On the Approximation of the Black Hole Shadow with a Simple Polar Curve}. \apj, 900, \href{https://iopscience.iop.org/article/10.3847/1538-4357/aba59a}{77}. 
%\href{https://arxiv.org/abs/2007.06732}{arXiv:2007.06732}.




\vspace{+.4cm}
\section{\sc Publications (\fcolorbox{red}{white}{Major} or normal contribution, \\$\ast$ indicates student-led) } 
\label{sec:publications}


(last updated: Oct 2026)

149. Shrestha, M., Kwok, L. A., Sand, D. J., ..., \textbf{J. R. Farah}, \textit{et al.} 2026. \textit{JWST Observations of SN 2024abup: First Detection of CO in a Broad-lined Type Ic Supernova and Constraints on r-process Nucleosynthesis}. \apjl, 1009, L53. \href{https://arxiv.org/abs/2606.28561}{arXiv:2606.28561}.

148. Macrie, C. W., Larison, C., Sears, H., ..., \textbf{J. R. Farah}, \textit{et al.} 2026. \textit{JWST Nebular Spectroscopy of SN 2023qov: Circumstellar Dust Emission in a Normal Type Ia Supernova}. \apj, 1009, 241. \href{https://arxiv.org/abs/2604.09777}{arXiv:2604.09777}.

147. Barna, B., Singh, M., Kwok, L. A., ..., \textbf{J. R. Farah}, \textit{et al.} 2026. \textit{Tomography analysis of the intermediate-luminosity Type Iax SN 2024pxl}. arXiv e-prints, \href{https://arxiv.org/abs/2609.28712}{arXiv:2609.28712}.

\begin{mdframed}[linecolor=red]*146. Duong, C., Myhre, F., \textbf{J. R. Farah}. 2026. \textit{Convolutional Neural Network for Extraction of $n = 1$ Photon Ring of Black Holes}. arXiv e-prints, \href{https://arxiv.org/abs/2609.25344}{arXiv:2609.25344}.\end{mdframed}

145. Dukiya, N., Misra, K., Pastorello, A., ..., \textbf{J. R. Farah}, \textit{et al.} 2026. \textit{Early NIR echo and the time variable mass loss history of Type IIn SN 2024kgi}. arXiv e-prints, \href{https://arxiv.org/abs/2609.25233}{arXiv:2609.25233}.

144. Song, C.-Y., Jiang, N., Wang, X., ..., \textbf{J. R. Farah}, \textit{et al.} 2026. \textit{SN 2025fhm: A central-engine powered Ic-BL supernova associated with X-ray transient EP250304a}. arXiv e-prints, \href{https://arxiv.org/abs/2609.08356}{arXiv:2609.08356}.

143. Dubey, M., Misra, K., Dukiya, N., ..., \textbf{J. R. Farah}, \textit{et al.} 2026. \textit{SN 2022xus: bridging the gap between Type IIP and IIL supernovae}. \mnras, 551, stag1427. \href{https://arxiv.org/abs/2607.25323}{arXiv:2607.25323}.

142. Hebbar, D., Roy, R., Nicholl, M., ..., \textbf{J. R. Farah}, \textit{et al.} 2026. \textit{SN 2019cqc: A Hydrogen-rich Superluminous Supernova Showing Electron-scattering Signatures}. \apj, 1008, 135. \href{https://arxiv.org/abs/2608.30309}{arXiv:2608.30309}.

141. Jiang, S.-Q., Xu, D., Lei, W.-H., ..., \textbf{J. R. Farah}, \textit{et al.} 2026. \textit{EP251023a: A fast X-ray transient featuring a magnetar-powered optical internal plateau followed by a steep decay}. \aap, 713, L5. \href{https://arxiv.org/abs/2606.24329}{arXiv:2606.24329}.

140. M{\"u}ller, H., von Fellenberg, S. D., Zeng, A.-L., ..., \textbf{J. R. Farah}, \textit{et al.} 2026. \textit{Dynamics and geometry of the inner sub-parsec-scale jet in 3C 279 observed with the Event Horizon Telescope}. \aap, 713, A274. \href{https://arxiv.org/abs/2607.16104}{arXiv:2607.16104}.

139. Principe, G., Algaba, J. C., Aviano, E., ..., \textbf{J. R. Farah}, \textit{et al.} 2026. \textit{Broadband multiwavelength properties of the archetypal blazar 3C 279 during the 2017 Event Horizon Telescope campaign}. \aap, 712, A224. \href{https://arxiv.org/abs/2606.25868}{arXiv:2606.25868}.

138. Aghakhanloo, M., Torres-Alb{\`a}, N., Smith, N., ..., \textbf{J. R. Farah}, \textit{et al.} 2026. \textit{AT 2016blu: Accretion-Powered Outbursts in a Luminous Blue Variable and Compact Object Binary}. arXiv e-prints, \href{https://arxiv.org/abs/2608.20469}{arXiv:2608.20469}.

137. Kwok, L. A., Blondin, S., Miller, A. A., ..., \textbf{J. R. Farah}, \textit{et al.} 2026. \textit{JWST Spectroscopy of Type Ia Supernova 2025rbs from Maximum Light to the Nebular Phase}. arXiv e-prints, \href{https://arxiv.org/abs/2608.10451}{arXiv:2608.10451}.

136. Dong, Y., Villar, V. A., Nugent, A., ..., \textbf{J. R. Farah}, \textit{et al.} 2026. \textit{Spectral Diversity in Type Ibn Supernovae and the Large Host Offset of SN 2024acyl}. \apj, 1007, 146. \href{https://arxiv.org/abs/2511.03926}{arXiv:2511.03926}.

135. Srinivasaragavan, G. P., Li, D., Hall, X. J., ..., \textbf{J. R. Farah}, \textit{et al.} 2026. \textit{EP250827b/SN 2025wkm: An X-Ray Flash-supernova Powered by a Central Engine and Circumstellar Interaction}. \apjl, 1006, L36. \href{https://arxiv.org/abs/2512.10239}{arXiv:2512.10239}.

134. Ransome, C. L., Sand, D. J., Bostroem, K. A., ..., \textbf{J. R. Farah}, \textit{et al.} 2026. \textit{The Eye of Sauron in SN 2025ngs: A Short-plateau Cousin of SN 1998S with Evidence for a Ring-like Circumstellar Medium}. \apj, 1006, 230. \href{https://arxiv.org/abs/2605.05380}{arXiv:2605.05380}.

133. Vasylyev, S. S., Hosseinzadeh, G., Dessart, L., ..., \textbf{J. R. Farah}, \textit{et al.} 2026. \textit{Ultraviolet to Infrared Spectroscopy of the Type Ibn SN 2023tsz Suggests a Lower-mass Progenitor}. arXiv e-prints, \href{https://arxiv.org/abs/2607.15342}{arXiv:2607.15342}.

132. Hsu, B., Smith, N., Bostroem, K. A., ..., \textbf{J. R. Farah}, \textit{et al.} 2026. \textit{Evidence for Asymmetric Ejecta and Circumstellar Material in SN 2023ixf Inferred from Extensive Nebular-phase Observations}. arXiv e-prints, \href{https://arxiv.org/abs/2607.02859}{arXiv:2607.02859}.

131. Hall, S., Kwok, L. A., Ravi, A. P., ..., \textbf{J. R. Farah}, \textit{et al.} 2026. \textit{JWST Observations of Calcium-Strong Transients: I. Complex Nebular He Emission in SN 2024uj}. arXiv e-prints, \href{https://arxiv.org/abs/2607.00111}{arXiv:2607.00111}.

130. Rosowsky, M., Valenti, S., Ravi, A. P., ..., \textbf{J. R. Farah}, \textit{et al.} 2026. \textit{SN 2023rve: A Type II Supernova with No Nebular Oxygen}. arXiv e-prints, \href{https://arxiv.org/abs/2606.30979}{arXiv:2606.30979}.

129. Ransome, C. L., Subrayan, B. M., Sand, D. J., ..., \textbf{J. R. Farah}, \textit{et al.} 2026. \textit{First Results from the LSST Shadow Survey: The Restless Luminous Blue Variable AT2017des in the Virgo-Cluster Galaxy, NGC4532}. arXiv e-prints, \href{https://arxiv.org/abs/2606.23784}{arXiv:2606.23784}.

128. Keinan, I., Arcavi, I., Howell, D. A., ..., \textbf{J. R. Farah}, \textit{et al.} 2026. \textit{Las Cumbres Observatory Gravitational-wave Follow-up in the Third and Fourth Observing Runs: Strengths and Weaknesses of a Rapid Response Galaxy Targeted Strategy}. \apj, 1004, 230. \href{https://arxiv.org/abs/2604.15129}{arXiv:2604.15129}.

127. Cai, J., Zeng, X., Zhou, W., ..., \textbf{J. R. Farah}, \textit{et al.} 2026. \textit{SN 2019dks: A 99aa-like Type Ia Supernova with a Low-velocity Gradient}. \apj, 1004, 98.

126. Bostroem, K. A., Valenti, S., Sand, D. J., ..., \textbf{J. R. Farah}, \textit{et al.} 2026. \textit{Late-time Hubble Space Telescope Ultraviolet Spectra of SN 2023ixf and SN 2024ggi Show Ongoing Interaction with Circumstellar Material}. \apj, 1004, 23. \href{https://arxiv.org/abs/2508.11756}{arXiv:2508.11756}.

125. Barna, B., B{\'a}nhidi, D., Szalai, T., ..., \textbf{J. R. Farah}, \textit{et al.} 2026. \textit{The extremely low-luminosity Type Iax SNe 2022ywf and 2023zgx}. \aap, 710, A72. \href{https://arxiv.org/abs/2602.23175}{arXiv:2602.23175}.

124. Hiramatsu, D., Berger, E., Tsuna, D., ..., \textbf{J. R. Farah}, \textit{et al.} 2026. \textit{The pair-instability origin of supernova 2023vbw}. arXiv e-prints, \href{https://arxiv.org/abs/2605.16487}{arXiv:2605.16487}.

123. Ravi, A. P., Kumar, S., Baer-Way, R., ..., \textbf{J. R. Farah}, \textit{et al.} 2026. \textit{The Double-peaked Calcium-strong SN 2025coe: Progenitor Constraints from Early Interaction and Ejecta Asymmetries}. \apj, 1003, 33. \href{https://arxiv.org/abs/2601.00415}{arXiv:2601.00415}.

122. Gerard, M., Andrews, J. E., Clayton, G. C., ..., \textbf{J. R. Farah}, \textit{et al.} 2026. \textit{SN 2024abfl: A Flat-plateau, Low-luminosity Type IIP Supernova with Early Circumstellar Matter Interaction}. \apj, 1002, 70. \href{https://arxiv.org/abs/2601.02638}{arXiv:2601.02638}.

121. Chen, L., Wang, X., Wu, Q., ..., \textbf{J. R. Farah}, \textit{et al.} 2026. \textit{SN 2024iss: A double-peaked Type IIb supernova with evidence of circumstellar interaction}. \aap, 710, A33. \href{https://arxiv.org/abs/2510.22997}{arXiv:2510.22997}.

120. K{\"o}nyves-T{\'o}th, R., Bodola, Z., Szalai, T., ..., \textbf{J. R. Farah}, \textit{et al.} 2026. \textit{Distance estimate to NGC 6951 from supernova siblings Type IIP SN 2020dpw and Type Ib SN 2021sjt}. \aap, 709, A68. \href{https://arxiv.org/abs/2603.24769}{arXiv:2603.24769}.

119. Aster, C., Inserra, C., Pastorello, A., ..., \textbf{J. R. Farah}, \textit{et al.} 2026. \textit{SN 2024abvb: a Type Ibn/Icn supernova with evidence of helium and an extreme offset from its host galaxy}. \mnras, 547, stag404. \href{https://arxiv.org/abs/2602.20775}{arXiv:2602.20775}.

118. Bernshteyn, V., Conroy, N. S., Baub{\"o}ck, M., ..., \textbf{J. R. Farah}, \textit{et al.} 2026. \textit{Ring Asymmetry and Spin in M87*}. \apj, 1000, 231. \href{https://arxiv.org/abs/2601.00394}{arXiv:2601.00394}.

117. Georgiev, B., Tiede, P., von Fellenberg, S. D., ..., \textbf{J. R. Farah}, \textit{et al.} 2026. \textit{Locating the missing large-scale emission in the jet of M87* with short EHT baselines}. \aap, 709, A26. \href{https://arxiv.org/abs/2601.13356}{arXiv:2601.13356}.

116. Albentosa-Ruiz, E., Washington, J. E., Marchili, N., ..., \textbf{J. R. Farah}, \textit{et al.} 2026. \textit{Full-polarization millimeter wavelength variability of Sagittarius A$^{*}$ during the 2018 EHT campaign}. \aap, 708, A179. \href{https://arxiv.org/abs/2604.10287}{arXiv:2604.10287}.

115. Singh, M., Kwok, L. A., Jha, S. W., ..., \textbf{J. R. Farah}, \textit{et al.} 2026. \textit{Photometry and Spectroscopy of SN 2024pxl: A Luminosity Link among Type Iax Supernovae}. \apj, 999, 227. \href{https://arxiv.org/abs/2505.02943}{arXiv:2505.02943}.

114. Dubey, M., Misra, K., Cs{\"o}rnyei, G., ..., \textbf{J. R. Farah}, \textit{et al.} 2026. \textit{SN 2023zcu: A Type IIP SN with Early Flash Features}. \apj, 999, 93. \href{https://arxiv.org/abs/2601.14830}{arXiv:2601.14830}.

113. Zimmerman, E. A., Gal-Yam, A., Groot, P. J., ..., \textbf{J. R. Farah}, \textit{et al.} 2026. \textit{A Faint Progenitor System for the Faint Supernova 2024vjm}. arXiv e-prints, \href{https://arxiv.org/abs/2602.09096}{arXiv:2602.09096}.

112. Kumar, H., Blanchard, P. K., Berger, E., ..., \textbf{J. R. Farah}, \textit{et al.} 2026. \textit{SN 2024afav: A Superluminous Supernova with Multiple Light-curve Bumps and Spectroscopic Signatures of Circumstellar Interaction}. \apjl, 998, L3. \href{https://arxiv.org/abs/2512.06067}{arXiv:2512.06067}.

111. Saurabh, M{\"u}ller, H., von Fellenberg, S. D., ..., \textbf{J. R. Farah}, \textit{et al.} 2026. \textit{Probing jet base emission of M87* with the 2021 Event Horizon Telescope observations}. \aap, 706, A27. \href{https://arxiv.org/abs/2512.08970}{arXiv:2512.08970}.

110. G{\'o}mez, J. L., Cho, I., Traianou, E., ..., \textbf{J. R. Farah}, \textit{et al.} 2026. \textit{Spatially resolved polarization swings in the supermassive binary black hole candidate OJ 287 with first Event Horizon Telescope observations}. \aap, 705, A23.

109. Event Horizon Telescope Collaboration, ..., \textbf{J. R. Farah}, \textit{et al.} 2025. \textit{Horizon-scale variability of M87* from 2017--2021 EHT observations}. \aap, 704, A91. \href{https://arxiv.org/abs/2509.24593}{arXiv:2509.24593}.

108. Park, S. H., Rho, J., Yoon, S.-C., ..., \textbf{J. R. Farah}, \textit{et al.} 2025. \textit{Near-infrared spectroscopy and detection of carbon monoxide in the Type II supernova SN 2023ixf}. \aap, 703, A227. \href{https://arxiv.org/abs/2507.11877}{arXiv:2507.11877}.

107. Kumar, H., Berger, E., Blanchard, P. K., ..., \textbf{J. R. Farah}, \textit{et al.} 2025. \textit{A Detection of Helium in the Bright Superluminous Supernova SN 2024rmj}. \apj, 992, 122. \href{https://arxiv.org/abs/2506.06417}{arXiv:2506.06417}.

106. Jacobson-Gal{\'a}n, W. V., Dessart, L., Davis, K. W., ..., \textbf{J. R. Farah}, \textit{et al.} 2025. \textit{Final Moments. III. Explosion Properties and Progenitor Constraints of CSM-interacting Type II Supernovae}. \apj, 992, 100. \href{https://arxiv.org/abs/2505.04698}{arXiv:2505.04698}.

105. Subrayan, B. M., Sand, D. J., Bostroem, K. A., ..., \textbf{J. R. Farah}, \textit{et al.} 2025. \textit{Early Shock Cooling Observations and Progenitor Constraints of Type IIb Supernova SN 2024uwq}. \apjl, 990, L68. \href{https://arxiv.org/abs/2505.02908}{arXiv:2505.02908}.

104. Shu, X., Yang, L., Yang, H., ..., \textbf{J. R. Farah}, \textit{et al.} 2025. \textit{EP241021a: A Months-duration X-Ray Transient with Luminous Optical and Radio Emission}. \apjl, 990, L29. \href{https://arxiv.org/abs/2505.07665}{arXiv:2505.07665}.

103. Hsu, B., Smith, N., Goldberg, J. A., ..., \textbf{J. R. Farah}, \textit{et al.} 2025. \textit{One Year of SN 2023ixf: Breaking through the Degenerate Parameter Space in Light-curve Models with Pulsating Progenitors}. \apj, 990, 148. \href{https://arxiv.org/abs/2408.07874}{arXiv:2408.07874}.

102. Pastorello, A., Reguitti, A., Tartaglia, L., ..., \textbf{J. R. Farah}, \textit{et al.} 2025. \textit{A long-lasting eruption heralds SN 2023ldh, a clone of SN 2009ip}. \aap, 701, A32. \href{https://arxiv.org/abs/2503.23123}{arXiv:2503.23123}.

101. Kwok, L. A., Singh, M., Jha, S. W., ..., \textbf{J. R. Farah}, \textit{et al.} 2025. \textit{JWST and Ground-based Observations of the Type Iax Supernovae SN 2024pxl and SN 2024vjm: Evidence for Weak Deflagration Explosions}. \apjl, 989, L33. \href{https://arxiv.org/abs/2505.02944}{arXiv:2505.02944}.

100. Gagliano, A., Villar, V. A., Matsumoto, T., ..., \textbf{J. R. Farah}, \textit{et al.} 2025. \textit{Evidence for an Instability-induced Binary Merger in the Double-peaked, Helium-rich Type IIn Supernova 2023zkd}. \apj, 989, 182. \href{https://arxiv.org/abs/2502.19469}{arXiv:2502.19469}.

99. {\'S}niegowska, M., Trakhtenbrot, B., Makrygianni, L., ..., \textbf{J. R. Farah}, \textit{et al.} 2025. \textit{AT 2019aalc: A Bowen Fluorescence Flare with a Precursor Flare in an Active Galactic Nucleus}. \apj, 989, 173. \href{https://arxiv.org/abs/2505.00083}{arXiv:2505.00083}.

98. Wang, Z.-Y., Pastorello, A., Cai, Y.-Z., ..., \textbf{J. R. Farah}, \textit{et al.} 2025. \textit{Massive stars exploding in a He-rich circumstellar medium: XI. Diverse evolution of five Ibn SNe 2020nxt, 2020taz, 2021bbv, 2023utc, and 2024aej}. \aap, 700, A156. \href{https://arxiv.org/abs/2506.15139}{arXiv:2506.15139}.

97. Sun, H., Li, W.-X., Liu, L.-D., ..., \textbf{J. R. Farah}, \textit{et al.} 2025. \textit{A fast X-ray transient from a weak relativistic jet associated with a type Ic-BL supernova}. Nature Astronomy, 9, 1073-1085. \href{https://arxiv.org/abs/2410.02315}{arXiv:2410.02315}.

96. Jiang, S.-Q., Xu, D., van Hoof, A. P. C., ..., \textbf{J. R. Farah}, \textit{et al.} 2025. \textit{EP240801a/XRF 240801B: An X-Ray Flash Detected by the Einstein Probe and the Implications of Its Multiband Afterglow}. \apjl, 988, L34. \href{https://arxiv.org/abs/2503.04306}{arXiv:2503.04306}.

95. Makrygianni, L., Arcavi, I., Newsome, M., ..., \textbf{J. R. Farah}, \textit{et al.} 2025. \textit{The Double Tidal Disruption Event AT 2022dbl Implies that at Least Some ``Standard'' Optical Tidal Disruption Events Are Partial Disruptions}. \apjl, 987, L20. \href{https://arxiv.org/abs/2505.16867}{arXiv:2505.16867}.

94. Kumar, H., Berger, E., Blanchard, P. K., ..., \textbf{J. R. Farah}, \textit{et al.} 2025. \textit{A Near-infrared Search for Helium in the Superluminous Supernova SN 2024ahr}. \apj, 987, 127. \href{https://arxiv.org/abs/2501.01485}{arXiv:2501.01485}.

93. Dahale, R., Cho, I., Moriyama, K., ..., \textbf{J. R. Farah}, \textit{et al.} 2025. \textit{Origin of the ring ellipticity in the black hole images of M87$^{*}$}. \aap, 699, A279. \href{https://arxiv.org/abs/2505.10333}{arXiv:2505.10333}.

92. Goddi, C., Carlos, D. F., Crew, G. B., ..., \textbf{J. R. Farah}, \textit{et al.} 2025. \textit{First polarization study of the M87 jet and active galactic nuclei at submillimeter wavelengths with ALMA}. \aap, 699, A265. \href{https://arxiv.org/abs/2505.10181}{arXiv:2505.10181}.

91. Palit, B., {\'S}niegowska, M., Markowitz, A., ..., \textbf{J. R. Farah}, \textit{et al.} 2025. \textit{Markarian 590: the AGN awakens}. \mnras, 540, L14-L20. \href{https://arxiv.org/abs/2501.07225}{arXiv:2501.07225}.

90. Onori, F., Nicholl, M., Ramsden, P., ..., \textbf{J. R. Farah}, \textit{et al.} 2025. \textit{The case of AT2022wtn: a tidal disruption event in an interacting galaxy}. \mnras, 540, 498-520. \href{https://arxiv.org/abs/2504.21686}{arXiv:2504.21686}.

89. Zeng, X., Hu, L., Zheng, S., ..., \textbf{J. R. Farah}, \textit{et al.} 2025. \textit{SN 2023ehl: A Normal Type Ia Supernova with High-velocity Features}. \apj, 986, 68.

88. Ma, X., Wang, X., Mo, J., ..., \textbf{J. R. Farah}, \textit{et al.} 2025. \textit{Supernovae at distances <40 Mpc: II. Supernova rate in the local Universe}. \aap, 698, A306. \href{https://arxiv.org/abs/2504.04507}{arXiv:2504.04507}.

87. Ma, X., Wang, X., Mo, J., ..., \textbf{J. R. Farah}, \textit{et al.} 2025. \textit{Supernovae at distances <40 Mpc: I. Catalogues and fractions of supernovae in a complete sample}. \aap, 698, A305. \href{https://arxiv.org/abs/2504.04393}{arXiv:2504.04393}.

86. Zeng, X., Luo, Y., Wang, X., ..., \textbf{J. R. Farah}, \textit{et al.} 2025. \textit{SN 2023xqm: A gradually fading Ia supernova exhibiting isolated high-velocity signatures}. \aap, 698, A70.

85. Liu, M., Zhang, Y., Wang, Y., ..., \textbf{J. R. Farah}, \textit{et al.} 2025. \textit{Detection of an Orphan X-Ray Flare from a Blazar Candidate EP240709a with the Einstein Probe}. \apj, 984, 5. \href{https://arxiv.org/abs/2412.18463}{arXiv:2412.18463}.

84. Li, W.-X., Zhu, Z.-P., Zou, X.-Z., ..., \textbf{J. R. Farah}, \textit{et al.} 2025. \textit{An extremely soft and weak fast X-ray transient associated with a luminous supernova}. arXiv e-prints, \href{https://arxiv.org/abs/2504.17034}{arXiv:2504.17034}.

83. Baer-Way, R., Chandra, P., Modjaz, M., ..., \textbf{J. R. Farah}, \textit{et al.} 2025. \textit{A Multiwavelength Autopsy of the Interacting Type IIn Supernova 2020ywx: Tracing Its Progenitor Mass-loss History for 100 Yr Before Death}. \apj, 983, 101. \href{https://arxiv.org/abs/2412.06914}{arXiv:2412.06914}.

82. Shrestha, M., DeSoto, S., Sand, D. J., ..., \textbf{J. R. Farah}, \textit{et al.} 2025. \textit{Spectropolarimetry of SN 2023ixf Reveals Both Circumstellar Material and an Aspherical Helium Core}. \apjl, 982, L32. \href{https://arxiv.org/abs/2410.08199}{arXiv:2410.08199}.

81. Ravi, A. P., Valenti, S., Dong, Y., ..., \textbf{J. R. Farah}, \textit{et al.} 2025. \textit{Luminous Type II Short-plateau SN 2023ufx: Asymmetric Explosion of a Partially Stripped Massive Progenitor}. \apj, 982, 12. \href{https://arxiv.org/abs/2411.02493}{arXiv:2411.02493}.

80. R{\"o}der, J., Wielgus, M., Lobanov, A. P., ..., \textbf{J. R. Farah}, \textit{et al.} 2025. \textit{A multifrequency study of sub-parsec jets with the Event Horizon Telescope}. \aap, 695, A233. \href{https://arxiv.org/abs/2501.05518}{arXiv:2501.05518}.

79. Andrews, J. E., Shrestha, M., Bostroem, K. A., ..., \textbf{J. R. Farah}, \textit{et al.} 2025. \textit{Asymmetries and Circumstellar Interaction in the Type II SN 2024bch}. \apj, 980, 37. \href{https://arxiv.org/abs/2411.02497}{arXiv:2411.02497}.

78. Yesmin, N., Pellegrino, C., Modjaz, M., ..., \textbf{J. R. Farah}, \textit{et al.} 2025. \textit{Spectral dataset of young type Ib supernovae and their time evolution}. \aap, 693, A307. \href{https://arxiv.org/abs/2409.04522}{arXiv:2409.04522}.

77. Event Horizon Telescope Collaboration, ..., \textbf{J. R. Farah}, \textit{et al.} 2025. \textit{The persistent shadow of the supermassive black hole of M87: II. Model comparisons and theoretical interpretations}. \aap, 693, A265.

76. Newsome, M., Arcavi, I., Howell, D. A., ..., \textbf{J. R. Farah}, \textit{et al.} 2024. \textit{Mapping the Inner 0.1 pc of a Supermassive Black Hole Environment with the Tidal Disruption Event and Extreme Coronal-line Emitter AT 2022upj}. \apj, 977, 258. \href{https://arxiv.org/abs/2406.11972}{arXiv:2406.11972}.

75. Dong, Y., Tsuna, D., Valenti, S., ..., \textbf{J. R. Farah}, \textit{et al.} 2024. \textit{SN2023fyq: A Type Ibn Supernova with Long-standing Precursor Activity Due to Binary Interaction}. \apj, 977, 254. \href{https://arxiv.org/abs/2405.04583}{arXiv:2405.04583}.

74. Pellegrino, C., Modjaz, M., Takei, Y., ..., \textbf{J. R. Farah}, \textit{et al.} 2024. \textit{The X-Ray Luminous Type Ibn SN 2022ablq: Estimates of Preexplosion Mass Loss and Constraints on Precursor Emission}. \apj, 977, 2. \href{https://arxiv.org/abs/2407.18291}{arXiv:2407.18291}.

73. Baczko, A.-K., Kadler, M., Ros, E., ..., \textbf{J. R. Farah}, \textit{et al.} 2024. \textit{The putative center in NGC 1052}. \aap, 692, A205. \href{https://arxiv.org/abs/2501.08685}{arXiv:2501.08685}.

72. Algaba, J. C., Balokovi{\'c}, M., Chandra, S., ..., \textbf{J. R. Farah}, \textit{et al.} 2024. \textit{Broadband multi-wavelength properties of M87 during the 2018 EHT campaign including a very high energy flaring episode}. \aap, 692, A140. \href{https://arxiv.org/abs/2404.17623}{arXiv:2404.17623}.

71. Kumar, H., Berger, E., Hiramatsu, D., ..., \textbf{J. R. Farah}, \textit{et al.} 2024. \textit{AT2023vto: An Exceptionally Luminous Helium Tidal Disruption Event from a Massive Star}. \apjl, 974, L36. \href{https://arxiv.org/abs/2408.01482}{arXiv:2408.01482}.

70. Dong, Y., Valenti, S., Ashall, C., ..., \textbf{J. R. Farah}, \textit{et al.} 2024. \textit{Characterizing the Rapid Hydrogen Disappearance in SN 2022crv: Evidence of a Continuum between Type Ib and IIb Supernova Properties}. \apj, 974, 316. \href{https://arxiv.org/abs/2309.09433}{arXiv:2309.09433}.

69. Szalai, T., K{\"o}nyves-T{\'o}th, R., Nagy, A. P., ..., \textbf{J. R. Farah}, \textit{et al.} 2024. \textit{The story of SN 2021aatd: A peculiar 1987A-like supernova with an early-phase luminosity excess}. \aap, 690, A17. \href{https://arxiv.org/abs/2406.02498}{arXiv:2406.02498}.

68. Bora, Z., K{\"o}nyves-T{\'o}th, R., Vink{\'o}, J., ..., \textbf{J. R. Farah}, \textit{et al.} 2024. \textit{Ejecta Masses in Type Ia Supernovae{\textemdash}Implications for the Progenitor and the Explosion Scenario}. \pasp, 136, 094201. \href{https://arxiv.org/abs/2408.11928}{arXiv:2408.11928}.

67. Shrestha, M., Bostroem, K. A., Sand, D. J., ..., \textbf{J. R. Farah}, \textit{et al.} 2024. \textit{Extended Shock Breakout and Early Circumstellar Interaction in SN 2024ggi}. \apjl, 972, L15. \href{https://arxiv.org/abs/2405.18490}{arXiv:2405.18490}.

66. Yadavalli, S. K., Villar, V. A., Izzo, L., ..., \textbf{J. R. Farah}, \textit{et al.} 2024. \textit{SN 2022oqm: A Bright and Multipeaked Calcium-rich Transient}. \apj, 972, 194. \href{https://arxiv.org/abs/2308.12991}{arXiv:2308.12991}.

65. Raymond, A. W., Doeleman, S. S., Asada, K., ..., \textbf{J. R. Farah}, \textit{et al.} 2024. \textit{First Very Long Baseline Interferometry Detections at 870 {\ensuremath{\mu}}m}. \aj, 168, 130. \href{https://arxiv.org/abs/2410.07453}{arXiv:2410.07453}.

64. Johnson, M. D., Akiyama, K., Baturin, R., ..., \textbf{J. R. Farah}, \textit{et al.} 2024. \textit{The Black Hole Explorer: motivation and vision}. Proc. SPIE, 13092, 130922D. \href{https://arxiv.org/abs/2406.12917}{arXiv:2406.12917}.

63. Jacobson-Gal{\'a}n, W. V., Dessart, L., Davis, K. W., ..., \textbf{J. R. Farah}, \textit{et al.} 2024. \textit{Final Moments. II. Observational Properties and Physical Modeling of Circumstellar-material-interacting Type II Supernovae}. \apj, 970, 189. \href{https://arxiv.org/abs/2403.02382}{arXiv:2403.02382}.

62. Faris, S., Arcavi, I., Makrygianni, L., ..., \textbf{J. R. Farah}, \textit{et al.} 2024. \textit{Light-curve Structure and H{\ensuremath{\alpha}} Line Formation in the Tidal Disruption Event AT 2019azh}. \apj, 969, 104. \href{https://arxiv.org/abs/2312.03842}{arXiv:2312.03842}.

61. Park, S. H., Rho, J., Yoon, S.-C., ..., \textbf{J. R. Farah}, \textit{et al.} 2024. \textit{Optical and Infrared Observation of a Type IIP Supernova 2021qqu}. ASP Conf. Ser., 536, 165.

60. Kwok, L. A., Siebert, M. R., Johansson, J., ..., \textbf{J. R. Farah}, \textit{et al.} 2024. \textit{Ground-based and JWST Observations of SN 2022pul. II. Evidence from Nebular Spectroscopy for a Violent Merger in a Peculiar Type Ia Supernova}. \apj, 966, 135. \href{https://arxiv.org/abs/2308.12450}{arXiv:2308.12450}.

59. Andrews, J. E., Pearson, J., Hosseinzadeh, G., ..., \textbf{J. R. Farah}, \textit{et al.} 2024. \textit{SN 2022jox: An Extraordinarily Ordinary Type II SN with Flash Spectroscopy}. \apj, 965, 85. \href{https://arxiv.org/abs/2310.16092}{arXiv:2310.16092}.

58. Event Horizon Telescope Collaboration, ..., \textbf{J. R. Farah}, \textit{et al.} 2024. \textit{First Sagittarius A* Event Horizon Telescope Results. VIII. Physical Interpretation of the Polarized Ring}. \apjl, 964, L26.

57. Event Horizon Telescope Collaboration, ..., \textbf{J. R. Farah}, \textit{et al.} 2024. \textit{First Sagittarius A* Event Horizon Telescope Results. VII. Polarization of the Ring}. \apjl, 964, L25.

56. Padilla Gonzalez, E., Howell, D. A., Terreran, G., ..., \textbf{J. R. Farah}, \textit{et al.} 2024. \textit{SN 2022joj: A Potential Double Detonation with a Thin Helium Shell}. \apj, 964, 196. \href{https://arxiv.org/abs/2308.06334}{arXiv:2308.06334}.

55. Hiramatsu, D., Matsumoto, T., Berger, E., ..., \textbf{J. R. Farah}, \textit{et al.} 2024. \textit{Multiple Peaks and a Long Precursor in the Type IIn Supernova 2021qqp: An Energetic Explosion in a Complex Circumstellar Environment}. \apj, 964, 181. \href{https://arxiv.org/abs/2305.11168}{arXiv:2305.11168}.

54. Shrestha, M., Pearson, J., Wyatt, S., ..., \textbf{J. R. Farah}, \textit{et al.} 2024. \textit{Evidence of Weak Circumstellar Medium Interaction in the Type II SN 2023axu}. \apj, 961, 247. \href{https://arxiv.org/abs/2310.00162}{arXiv:2310.00162}.

53. Newsome, M., Arcavi, I., Howell, D. A., ..., \textbf{J. R. Farah}, \textit{et al.} 2024. \textit{Probing the Subparsec Dust of a Supermassive Black Hole with the Tidal Disruption Event AT 2020mot}. \apj, 961, 239. \href{https://arxiv.org/abs/2305.03767}{arXiv:2305.03767}.

52. Paraschos, G. F., Kim, J.-Y., Wielgus, M., ..., \textbf{J. R. Farah}, \textit{et al.} 2024. \textit{Ordered magnetic fields around the 3C 84 central black hole}. \aap, 682, L3. \href{https://arxiv.org/abs/2402.00927}{arXiv:2402.00927}.

51. Siebert, M. R., Kwok, L. A., Johansson, J., ..., \textbf{J. R. Farah}, \textit{et al.} 2024. \textit{Ground-based and JWST Observations of SN 2022pul. I. Unusual Signatures of Carbon, Oxygen, and Circumstellar Interaction in a Peculiar Type Ia Supernova}. \apj, 960, 88. \href{https://arxiv.org/abs/2308.12449}{arXiv:2308.12449}.

50. Pearson, J., Sand, D. J., Lundqvist, P., ..., \textbf{J. R. Farah}, \textit{et al.} 2024. \textit{Strong Carbon Features and a Red Early Color in the Underluminous Type Ia SN 2022xkq}. \apj, 960, 29. \href{https://arxiv.org/abs/2309.10054}{arXiv:2309.10054}.

49. Event Horizon Telescope Collaboration, ..., \textbf{J. R. Farah}, \textit{et al.} 2024. \textit{The persistent shadow of the supermassive black hole of M 87. I. Observations, calibration, imaging, and analysis}. \aap, 681, A79.

48. Torne, P., Liu, K., Eatough, R. P., ..., \textbf{J. R. Farah}, \textit{et al.} 2023. \textit{A Search for Pulsars around Sgr A* in the First Event Horizon Telescope Data Set}. \apj, 959, 14. \href{https://arxiv.org/abs/2308.15381}{arXiv:2308.15381}.

47. Roelofs, F., Johnson, M. D., Chael, A., ..., \textbf{J. R. Farah}, \textit{et al.} 2023. \textit{Polarimetric Geometric Modeling for mm-VLBI Observations of Black Holes}. \apjl, 957, L21. \href{https://arxiv.org/abs/2311.10695}{arXiv:2311.10695}.

46. Event Horizon Telescope Collaboration, ..., \textbf{J. R. Farah}, \textit{et al.} 2023. \textit{First M87 Event Horizon Telescope Results. IX. Detection of Near-horizon Circular Polarization}. \apjl, 957, L20. \href{https://arxiv.org/abs/2311.10976}{arXiv:2311.10976}.

45. Bostroem, K. A., Pearson, J., Shrestha, M., ..., \textbf{J. R. Farah}, \textit{et al.} 2023. \textit{Early Spectroscopy and Dense Circumstellar Medium Interaction in SN 2023ixf}. \apjl, 956, L5. \href{https://arxiv.org/abs/2306.10119}{arXiv:2306.10119}.

44. Hiramatsu, D., Tsuna, D., Berger, E., ..., \textbf{J. R. Farah}, \textit{et al.} 2023. \textit{From Discovery to the First Month of the Type II Supernova 2023ixf: High and Variable Mass Loss in the Final Year before Explosion}. \apjl, 955, L8. \href{https://arxiv.org/abs/2307.03165}{arXiv:2307.03165}.

43. Bostroem, K. A., Dessart, L., Hillier, D. J., ..., \textbf{J. R. Farah}, \textit{et al.} 2023. \textit{SN 2022acko: The First Early Far-ultraviolet Spectra of a Type IIP Supernova}. \apjl, 953, L18. \href{https://arxiv.org/abs/2305.01654}{arXiv:2305.01654}.

\begin{mdframed}[linecolor=red]
42. Hosseinzadeh, G., \textbf{J. R. Farah}, Shrestha, M., \textit{et al.} 2023. \textit{Shock Cooling and Possible Precursor Emission in the Early Light Curve of the Type II SN 2023ixf}. \apjl, 953, L16. \href{https://arxiv.org/abs/2306.06097}{arXiv:2306.06097}.

\end{mdframed}

41. Hosseinzadeh, G., Sand, D. J., Sarbadhicary, S. K., ..., \textbf{J. R. Farah}, \textit{et al.} 2023. \textit{The Early Light Curve of SN 2023bee: Constraining Type Ia Supernova Progenitors the Apian Way}. \apjl, 953, L15. \href{https://arxiv.org/abs/2305.03071}{arXiv:2305.03071}.

40. Prather, B. S., Dexter, J., Moscibrodzka, M., ..., \textbf{J. R. Farah}, \textit{et al.} 2023. \textit{Comparison of Polarized Radiative Transfer Codes Used by the EHT Collaboration}. \apj, 950, 35. \href{https://arxiv.org/abs/2303.12004}{arXiv:2303.12004}.

39. Shrestha, M., Sand, D. J., Alexander, K. D., ..., \textbf{J. R. Farah}, \textit{et al.} 2023. \textit{Limit on Supernova Emission in the Brightest Gamma-Ray Burst, GRB 221009A}. \apjl, 946, L25. \href{https://arxiv.org/abs/2302.03829}{arXiv:2302.03829}.

38. Ben-Ami, T., Arcavi, I., Newsome, M., \textbf{J. R. Farah}, \textit{et al.} 2023. \textit{The Type Ibn Supernova 2019kbj: Indications for Diversity in Type Ibn Supernova Progenitors}. \apj, 946, 30. \href{https://arxiv.org/abs/2212.03407}{arXiv:2212.03407}.

37. Jorstad, S., Wielgus, M., Lico, R., ..., \textbf{J. R. Farah}, \textit{et al.} 2023. \textit{The Event Horizon Telescope Image of the Quasar NRAO 530}. \apj, 943, 170. \href{https://arxiv.org/abs/2302.04622}{arXiv:2302.04622}.

\begin{mdframed}[linecolor=red]
36. Wielgus, M., Moscibrodzka, M., Vos, J., ..., \textbf{J. R. Farah}, \textit{et al.} 2022. \textit{Orbital motion near Sagittarius A$^{*}$ . Constraints from polarimetric ALMA observations}. \aap, 665, L6. \href{https://arxiv.org/abs/2209.09926}{arXiv:2209.09926}.

\end{mdframed}

35. Broderick, A. E., Pesce, D. W., Gold, R., ..., \textbf{J. R. Farah}, \textit{et al.} 2022. \textit{The Photon Ring in M87*}. \apj, 935, 61. \href{https://arxiv.org/abs/2208.09004}{arXiv:2208.09004}.

34. Issaoun, S., Wielgus, M., Jorstad, S., ..., \textbf{J. R. Farah}, \textit{et al.} 2022. \textit{Resolving the Inner Parsec of the Blazar J1924--2914 with the Event Horizon Telescope}. \apj, 934, 145. \href{https://arxiv.org/abs/2208.01662}{arXiv:2208.01662}.

\begin{mdframed}[linecolor=red]
33. {\renewcommand{\thefootnote}{\ensuremath{S}}\renewcommand{\thempfootnote}{\ensuremath{S}}\footnote{This is the first Galactic Center results series from the Event Horizon Telescope, presenting the second image of a black hole and the first horizon-scale images of Sgr A$^\ast$.}}Broderick, A. E., Gold, R., Georgiev, B., ..., \textbf{J. R. Farah}, \textit{et al.} 2022. \textit{Characterizing and Mitigating Intraday Variability: Reconstructing Source Structure in Accreting Black Holes with mm-VLBI}. \apjl, 930, L21.

32. $^S$Georgiev, B., Pesce, D. W., Broderick, A. E., ..., \textbf{J. R. Farah}, \textit{et al.} 2022. \textit{A Universal Power-law Prescription for Variability from Synthetic Images of Black Hole Accretion Flows}. \apjl, 930, L20.

31. $^S$Event Horizon Telescope Collaboration, ..., \textbf{J. R. Farah}, \textit{et al.} 2022. \textit{First Sagittarius A* Event Horizon Telescope Results. VI. Testing the Black Hole Metric}. \apjl, 930, L17. \href{https://arxiv.org/abs/2311.09484}{arXiv:2311.09484}.

30. $^S$Event Horizon Telescope Collaboration, ..., \textbf{J. R. Farah}, \textit{et al.} 2022. \textit{First Sagittarius A* Event Horizon Telescope Results. V. Testing Astrophysical Models of the Galactic Center Black Hole}. \apjl, 930, L16. \href{https://arxiv.org/abs/2311.09478}{arXiv:2311.09478}.

29. $^S$Event Horizon Telescope Collaboration, ..., \textbf{J. R. Farah}, \textit{et al.} 2022. \textit{First Sagittarius A* Event Horizon Telescope Results. IV. Variability, Morphology, and Black Hole Mass}. \apjl, 930, L15. \href{https://arxiv.org/abs/2311.08697}{arXiv:2311.08697}.

28. $^S$Event Horizon Telescope Collaboration, ..., \textbf{J. R. Farah}, \textit{et al.} 2022. \textit{First Sagittarius A* Event Horizon Telescope Results. III. Imaging of the Galactic Center Supermassive Black Hole}. \apjl, 930, L14. \href{https://arxiv.org/abs/2311.09479}{arXiv:2311.09479}.

27. $^S$Event Horizon Telescope Collaboration, ..., \textbf{J. R. Farah}, \textit{et al.} 2022. \textit{First Sagittarius A* Event Horizon Telescope Results. II. EHT and Multiwavelength Observations, Data Processing, and Calibration}. \apjl, 930, L13. \href{https://arxiv.org/abs/2311.08679}{arXiv:2311.08679}.

26. $^S$Event Horizon Telescope Collaboration, ..., \textbf{J. R. Farah}, \textit{et al.} 2022. \textit{First Sagittarius A* Event Horizon Telescope Results. I. The Shadow of the Supermassive Black Hole in the Center of the Milky Way}. \apjl, 930, L12. \href{https://arxiv.org/abs/2311.08680}{arXiv:2311.08680}.

\end{mdframed}

25. Wielgus, M., Marchili, N., Mart{\'\i}-Vidal, I., ..., \textbf{J. R. Farah}, \textit{et al.} 2022. \textit{Millimeter Light Curves of Sagittarius A* Observed during the 2017 Event Horizon Telescope Campaign}. \apjl, 930, L19. \href{https://arxiv.org/abs/2207.06829}{arXiv:2207.06829}.

24. Satapathy, K., Psaltis, D., {\"O}zel, F., ..., \textbf{J. R. Farah}, \textit{et al.} 2022. \textit{The Variability of the Black Hole Image in M87 at the Dynamical Timescale}. \apj, 925, 13. \href{https://arxiv.org/abs/2111.01317}{arXiv:2111.01317}.

23. Janssen, M., Falcke, H., Kadler, M., ..., \textbf{J. R. Farah}, \textit{et al.} 2021. \textit{Event Horizon Telescope observations of the jet launching and collimation in Centaurus A}. Nature Astronomy, 5, 1017-1028. \href{https://arxiv.org/abs/2111.03356}{arXiv:2111.03356}.

22. Kocherlakota, P., Rezzolla, L., Falcke, H., ..., \textbf{J. R. Farah}, \textit{et al.} 2021. \textit{Constraints on black-hole charges with the 2017 EHT observations of M87*}. \prd, 103, 104047. \href{https://arxiv.org/abs/2105.09343}{arXiv:2105.09343}.

21. Narayan, R., Palumbo, D. C. M., Johnson, M. D., ..., \textbf{J. R. Farah}, \textit{et al.} 2021. \textit{The Polarized Image of a Synchrotron-emitting Ring of Gas Orbiting a Black Hole}. \apj, 912, 35. \href{https://arxiv.org/abs/2105.01804}{arXiv:2105.01804}.

20. EHT MWL Science Working Group, ..., \textbf{J. R. Farah}, \textit{et al.} 2021. \textit{Broadband Multi-wavelength Properties of M87 during the 2017 Event Horizon Telescope Campaign}. \apjl, 911, L11. \href{https://arxiv.org/abs/2104.06855}{arXiv:2104.06855}.

19. Raymond, A. W., Palumbo, D., Paine, S. N., ..., \textbf{J. R. Farah}, \textit{et al.} 2021. \textit{Evaluation of New Submillimeter VLBI Sites for the Event Horizon Telescope}. \apjs, 253, 5. \href{https://arxiv.org/abs/2102.05482}{arXiv:2102.05482}.

18. Goddi, C., Mart{\'\i}-Vidal, I., Messias, H., ..., \textbf{J. R. Farah}, \textit{et al.} 2021. \textit{Polarimetric Properties of Event Horizon Telescope Targets from ALMA}. \apjl, 910, L14. \href{https://arxiv.org/abs/2105.02272}{arXiv:2105.02272}.

17. Event Horizon Telescope Collaboration, ..., \textbf{J. R. Farah}, \textit{et al.} 2021. \textit{First M87 Event Horizon Telescope Results. VIII. Magnetic Field Structure near The Event Horizon}. \apjl, 910, L13. \href{https://arxiv.org/abs/2105.01173}{arXiv:2105.01173}.

16. Event Horizon Telescope Collaboration, ..., \textbf{J. R. Farah}, \textit{et al.} 2021. \textit{First M87 Event Horizon Telescope Results. VII. Polarization of the Ring}. \apjl, 910, L12. \href{https://arxiv.org/abs/2105.01169}{arXiv:2105.01169}.

15. Kim, J.-Y., Krichbaum, T. P., Broderick, A. E., ..., \textbf{J. R. Farah}, \textit{et al.} 2020. \textit{Event Horizon Telescope imaging of the archetypal blazar 3C 279 at an extreme 20 microarcsecond resolution}. \aap, 640, A69.

14. Johnson, M. D., Lupsasca, A., Strominger, A., ..., \textbf{J. R. Farah}, \textit{et al.} 2020. \textit{Universal interferometric signatures of a black hole's photon ring}. Science Advances, 6, eaaz1310. \href{https://arxiv.org/abs/1907.04329}{arXiv:1907.04329}.

13. Fabbiano, G., Elvis, M., Accomazzi, A., ..., \textbf{J. R. Farah}, \textit{et al.} 2019. \textit{Increasing the Discovery Space in Astrophysics: The Exploration Question for Compact Objects}. \baas, 51, 89.

12. Fabbiano, G., Elvis, M., Accomazzi, A., ..., \textbf{J. R. Farah}, \textit{et al.} 2019. \textit{Increasing the Discovery Space in Astrophysics: The Exploration Question for Cosmology}. \baas, 51, 88.

11. Fabbiano, G., Elvis, M., Accomazzi, A., ..., \textbf{J. R. Farah}, \textit{et al.} 2019. \textit{Increasing the Discovery Space in Astrophysics: The Exploration Question for Galaxy Evolution}. \baas, 51, 87.

10. Fabbiano, G., Elvis, M., Accomazzi, A., ..., \textbf{J. R. Farah}, \textit{et al.} 2019. \textit{Increasing the Discovery Space in Astrophysics: The Exploration Question for Planetary Systems}. \baas, 51, 86.

9. Fabbiano, G., Elvis, M., Accomazzi, A., ..., \textbf{J. R. Farah}, \textit{et al.} 2019. \textit{Increasing the Discovery Space in Astrophysics: The Exploration Question for Stars and Stellar Evolution}. \baas, 51, 85.

8. Fabbiano, G., Elvis, M., Accomazzi, A., ..., \textbf{J. R. Farah}, \textit{et al.} 2019. \textit{Increasing the Discovery Space in Astrophysics: The Exploration Question for Resolved Stellar Populations}. \baas, 51, 84.

\begin{mdframed}[linecolor=red]
7. {\renewcommand{\thefootnote}{\ensuremath{\dagger}}\renewcommand{\thempfootnote}{\ensuremath{\dagger}}\footnote{These six papers are the first results papers from the Event Horizon Telescope, presenting the first image of a black hole.}}Event Horizon Telescope Collaboration, ..., \textbf{J. R. Farah}, \textit{et al.} 2019. \textit{First M87 Event Horizon Telescope Results. VI. The Shadow and Mass of the Central Black Hole}. \apjl, 875, L6. \href{https://arxiv.org/abs/1906.11243}{arXiv:1906.11243}.

6. $^\dagger$Event Horizon Telescope Collaboration, ..., \textbf{J. R. Farah}, \textit{et al.} 2019. \textit{First M87 Event Horizon Telescope Results. V. Physical Origin of the Asymmetric Ring}. \apjl, 875, L5. \href{https://arxiv.org/abs/1906.11242}{arXiv:1906.11242}.

5. $^\dagger$Event Horizon Telescope Collaboration, ..., \textbf{J. R. Farah}, \textit{et al.} 2019. \textit{First M87 Event Horizon Telescope Results. IV. Imaging the Central Supermassive Black Hole}. \apjl, 875, L4. \href{https://arxiv.org/abs/1906.11241}{arXiv:1906.11241}.

4. $^\dagger$Event Horizon Telescope Collaboration, ..., \textbf{J. R. Farah}, \textit{et al.} 2019. \textit{First M87 Event Horizon Telescope Results. III. Data Processing and Calibration}. \apjl, 875, L3. \href{https://arxiv.org/abs/1906.11240}{arXiv:1906.11240}.

3. $^\dagger$Event Horizon Telescope Collaboration, ..., \textbf{J. R. Farah}, \textit{et al.} 2019. \textit{First M87 Event Horizon Telescope Results. II. Array and Instrumentation}. \apjl, 875, L2. \href{https://arxiv.org/abs/1906.11239}{arXiv:1906.11239}.

2. $^\dagger$Event Horizon Telescope Collaboration, ..., \textbf{J. R. Farah}, \textit{et al.} 2019. \textit{First M87 Event Horizon Telescope Results. I. The Shadow of the Supermassive Black Hole}. \apjl, 875, L1. \href{https://arxiv.org/abs/1906.11238}{arXiv:1906.11238}.

\end{mdframed}

1. Fabbiano, G., Elvis, M., Accomazzi, A., ..., \textbf{J. R. Farah}, \textit{et al.} 2019. \textit{Increasing the Discovery Space in Astrophysics - A Collation of Six Submitted White Papers}. arXiv e-prints, \href{https://arxiv.org/abs/1903.06634}{arXiv:1903.06634}.


% \vspace{+.4cm}
% \section{\sc Conference publications}
% (last updated: Aug 2025)

% 4. M. D. Johnson, K. Akiyama, R. Baturin, ... \textbf{J. Farah} et al. (2024). \textit{The Black Hole Explorer: motivation and vision}. In: Space Telescopes and Instrumentation 2024: Optical, Infrared, and Millimeter Wave (SPIE Conf. Ser. 13092), 130922D.

% 3. Haworth, K., Johnson, M.~D., Pesce, D.~W., ... \textbf{J. R. Farah}  et al.\ 2019. \textit{Studying black holes on horizon scales with space-VLBI}. arXiv:1909.01405 

% 2. Blackburn, L., Doeleman, S., Dexter, J., ... \textbf{J. R. Farah} et al.\ 2019. \textit{Studying Black Holes on Horizon Scales with VLBI Ground Arrays}. arXiv:1909.01411 

% 1. Fabbiano, E., Berriman, B., Bose, C., ... \textbf{J. R. Farah}, et al. 2019. \textit{Increasing the Discovery Space in Astrophysics--A Collation of Six Submitted White Papers.} 


\vspace{+.4cm}
\section{\sc Technical publications}
(last updated: Aug 2025)

5. Chael A., Bouman, K., Johnson  M., Wielgus M., Blackburn L., Chan C., \textbf{Farah J.}, Palumbo D., Pesce D., et al. 2022. \textit{eht-imaging: v1.2.4: Imaging interferometric data with regularized maximum likelihood}. Zenodo (software) publication.

4. Chael A., Bouman, K., Johnson  M., Wielgus M., Blackburn L., Chan C., \textbf{Farah J.}, Palumbo D., Pesce D. 2019. \textit{eht-imaging: v1.1.0: Imaging interferometric data with regularized maximum likelihood}. Zenodo (software) publication.

3. \textbf{Farah J}, Felt N, Franklin M, Giorimini P, Rogan C, Tuna A, Wang A. 2017. \textit{All The Sparks We Cannot See}. ATLAS Internal Note.

2. \textbf{Farah J}, Felt N, Frank, Giorimini P, Rogan C, Tuna A, Wang A. 2017. \textit{Test of the Micromegas Trigger Processor with Cosmic Ray Muons}. ATLAS Internal Note.

1. Felt N, Franklin M, Giorimini P, Rogan C, Tuna A, Wang A, \textbf{Farah J}. 2017. \textit{Performance of a Micromegas octuplet after removing the major cause of MMFE8 noise}. ATLAS Internal Note.

\vspace{+.4cm}
\section{\sc Classifications, telegrams, and circulars}
(last updated: Aug 2025)

\vspace{-2mm}
First or co-author on 90+ Transient Name Server classification reports, AstroNotes, and GCN circulars, 2022--present, as part of the Global Supernova Project.

\vspace{+.2cm}
\section{\sc Software Development}

\texttt{eht-imaging}. Contributor. {\color{blue}\href{https://github.com/achael/eht-imaging}{Repository.}}

\vspace*{-2mm}
The Supernova Exchange (SNEx). Contributor. {\color{blue}\href{https://supernova.exchange}{Live site link.}}



% %% using icarus from ADS
% \bibitem[Farah et al.(2025)]{2025TNSCR3013....1F} \textbf{Farah, J.}, Andrews, M., Wynn, K., Howell, D.~A., McCully, C.\ 2025.\ Global SN Project Transient Classification Report for 2025-07-31.\ Transient Name Server Classification Report 2025-3013.
% \bibitem[Andrews et al.(2025)]{2025TNSCR2699....1A} Andrews, M., \textbf{Farah, J.}, Wynn, K., Bosteroem, A., Howell, D.~A., McCully, C.\ 2025.\ Global SN Project Transient Classification Report for 2025-07-14.\ Transient Name Server Classification Report 2025-2699.
% \bibitem[Li et al.(2025)]{2025GCN.40975....1L} Li, W.~X. and 8 colleagues 2025.\ EP250704a/GRB 250704B: GSP observations of the optical counterpart.\ GRB Coordinates Network, Circular Service, No. 40975 40975.
% \bibitem[Farah et al.(2025)]{2025TNSCR2162....1F} \textbf{Farah, J.}, Andrew, M., Howell, D.~A., McCully, C., Ni, Y.~Q., Wynn, K.\ 2025.\ Global SN Project Transient Classification Report for 2025-06-09.\ Transient Name Server Classification Report 2025-2162.
% \bibitem[Ni et al.(2025)]{2025TNSCR2041....1N} Ni, C., Andrews, M., \textbf{Farah, J.}, Wynn, K., Howell, D.~A., McCully, C.\ 2025.\ Global SN Project Transient Classification Report for 2025-05-30.\ Transient Name Server Classification Report 2025-2041.
% \bibitem[Ni et al.(2025)]{2025TNSCR2027....1N} Ni, C., Andrews, M., \textbf{Farah, J.}, Wynn, K., Howell, D.~A., McCully, C.\ 2025.\ Global SN Project Transient Classification Report for 2025-05-29.\ Transient Name Server Classification Report 2025-2027.
% \bibitem[Andrews et al.(2025)]{2025TNSCR1947....1A} Andrews, M., \textbf{Farah, J.}, Ni, C., Wynn, K., Howell, D.~A., McCully, C.\ 2025.\ Global SN Project Transient Classification Report for 2025-05-24.\ Transient Name Server Classification Report 2025-1947.
% \bibitem[Ni et al.(2025)]{2025TNSCR1733....1N} Ni, C., Andrews, M., \textbf{Farah, J.}, Wynn, K., Howell, D.~A., McCully, C.\ 2025.\ Global SN Project Transient Classification Report for 2025-05-09.\ Transient Name Server Classification Report 2025-1733.
% \bibitem[Baldini et al.(2025)]{2025TNSCR1720....1B} Baldini, P. and 7 colleagues 2025.\ Transient Classification Report for 2025-05-08.\ Transient Name Server Classification Report 2025-1720.
% \bibitem[Baldini et al.(2025)]{2025TNSAN.142....1B} Baldini, P. and 7 colleagues 2025.\ Spectroscopic classification of the re-brightening associated with AT 2021gqu as a SNIa-pec..\ Transient Name Server AstroNote 142.
% \bibitem[Andrews et al.(2025)]{2025TNSCR1539....1A} Andrews, M., \textbf{Farah, J.}, Ni, C., Wynn, K., Howell, D.~A., McCully, C.\ 2025.\ Global SN Project Transient Classification Report for 2025-04-25.\ Transient Name Server Classification Report 2025-1539.
% \bibitem[Andrews et al.(2025)]{2025TNSCR1531....1A} Andrews, M., \textbf{Farah, J.}, Ni, C., Wynn, K., Howell, D.~A., McCully, C.\ 2025.\ Global SN Project Transient Classification Report for 2025-04-24.\ Transient Name Server Classification Report 2025-1531.
% \bibitem[Newsome et al.(2025)]{2025TNSCR1351....1N} Newsome, M. and 8 colleagues 2025.\ Transient Classification Report for 2025-04-10.\ Transient Name Server Classification Report 2025-1351.
% \bibitem[Li et al.(2025)]{2025GCN.40157....1L} Li, W.~X. and 8 colleagues 2025.\ EP250416a: GSP optical upper limit.\ GRB Coordinates Network, Circular Service, No. 40157 40157.
% \bibitem[Andrews et al.(2025)]{2025TNSCR1076....1A} Andrews, M., \textbf{Farah, J.}, Howell, D.~A., McCully, C.\ 2025.\ Transient Classification Report for 2025-03-21.\ Transient Name Server Classification Report 2025-1076.
% \bibitem[Andrews et al.(2025)]{2025TNSCR1033....1A} Andrews, M., \textbf{Farah, J.}, Howell, D.~A., McCully, C.\ 2025.\ Transient Classification Report for 2025-03-18.\ Transient Name Server Classification Report 2025-1033.
% \bibitem[Andrews et al.(2025)]{2025TNSCR.991....1A} Andrews, M., \textbf{Farah, J.}, Howell, D.~A., McCully, C.\ 2025.\ Transient Classification Report for 2025-03-14.\ Transient Name Server Classification Report 2025-991.
% \bibitem[Andrews et al.(2025)]{2025TNSCR.977....1A} Andrews, M., \textbf{Farah, J.}, Howell, D.~A., McCully, C.\ 2025.\ Transient Classification Report for 2025-03-13.\ Transient Name Server Classification Report 2025-977.
% \bibitem[Farah et al.(2025)]{2025TNSCR.922....1F} \textbf{Farah, J.}, Andrews, M., Ni, C., Howell, D.~A., McCully, C.\ 2025.\ Transient Classification Report for 2025-03-08.\ Transient Name Server Classification Report 2025-922.
% \bibitem[Andrews et al.(2025)]{2025TNSCR.914....1A} Andrews, M. and 6 colleagues 2025.\ Transient Classification Report for 2025-03-07.\ Transient Name Server Classification Report 2025-914.
% \bibitem[Shrestha et al.(2025)]{2025GCN.39896....1S} Shrestha, M. and 13 colleagues 2025.\ GRB 250327B: Las Cumbres optical counterpart detection.\ GRB Coordinates Network, Circular Service, No. 39896 39896.
% \bibitem[Li et al.(2025)]{2025GCN.39627....1L} Li, W.~X. and 8 colleagues 2025.\ EP250308a: GSP optical upper limit.\ GRB Coordinates Network, Circular Service, No. 39627 39627.
% \bibitem[Andrews et al.(2025)]{2025TNSCR.802....1A} Andrews, M., Hiramatsu, D., \textbf{Farah, J.}, Howell, D.~A., McCully, C.\ 2025.\ Transient Classification Report for 2025-02-26.\ Transient Name Server Classification Report 2025-802.
% \bibitem[Farah et al.(2025)]{2025TNSCR.530....1F} \textbf{Farah, J.}, Andrews, M., Ni, C., Howell, D.~A., McCully, C., Wynn, K.\ 2025.\ Transient Classification Report for 2025-02-07.\ Transient Name Server Classification Report 2025-530.
% \bibitem[Li et al.(2025)]{2025GCN.39489....1L} Li, W.~X. and 8 colleagues 2025.\ EP250226a/GRB 250226A: GSP detects optical counterpart.\ GRB Coordinates Network, Circular Service, No. 39489 39489.
% \bibitem[Li et al.(2025)]{2025GCN.39477....1L} Li, W.~X. and 8 colleagues 2025.\ EP250225a: GSP optical upper limit.\ GRB Coordinates Network, Circular Service, No. 39477 39477.
% \bibitem[Faris et al.(2025)]{2025TNSCR..25....1F} Faris, S. and 6 colleagues 2025.\ StarDestroyers Transient Classification Report for 2025-01-02.\ Transient Name Server Classification Report 2025-25.
% \bibitem[Andrews et al.(2024)]{2024TNSCR4593....1A} Andrews, M., \textbf{Farah, J.}, Ni, C., Howell, D.~A., McCully, C.\ 2024.\ Global SN Project Transient Classification Report for 2024-11-22.\ Transient Name Server Classification Report 2024-4593.
% \bibitem[Newsome et al.(2024)]{2024TNSCR4486....1N} Newsome, M. and 6 colleagues 2024.\ Transient Classification Report for 2024-11-13.\ Transient Name Server Classification Report 2024-4486.
% \bibitem[Newsome et al.(2024)]{2024TNSCR4446....1N} Newsome, M. and 6 colleagues 2024.\ Transient Classification Report for 2024-11-11.\ Transient Name Server Classification Report 2024-4446.
% \bibitem[Faris et al.(2024)]{2024TNSCR4428....1F} Faris, S. and 6 colleagues 2024.\ StarDestroyers Transient Classification Report for 2024-11-10.\ Transient Name Server Classification Report 2024-4428.
% \bibitem[Newsome et al.(2024)]{2024TNSCR4337....1N} Newsome, M. and 6 colleagues 2024.\ Transient Classification Report for 2024-11-05.\ Transient Name Server Classification Report 2024-4337.
% \bibitem[Li et al.(2024)]{2024GCN.38338....1L} Li, W.~X. and 8 colleagues 2024.\ EP241126a: GSP detects optical counterpart.\ GRB Coordinates Network, Circular Service, No. 38338 38338.
% \bibitem[Li et al.(2024)]{2024GCN.38127....1L} Li, W.~X. and 8 colleagues 2024.\ EP241107a: GSP detects optical counterpart.\ GRB Coordinates Network, Circular Service, No. 38127 38127.
% \bibitem[Faris et al.(2024)]{2024TNSCR4005....1F} Faris, S. and 6 colleagues 2024.\ StarDestroyers Transient Classification Report for 2024-10-14.\ Transient Name Server Classification Report 2024-4005.
% \bibitem[Shrestha et al.(2024)]{2024GCN.37913....1S} Shrestha, M. and 11 colleagues 2024.\ GRB 241026A: Las Cumbres optical detection.\ GRB Coordinates Network, Circular Service, No. 37913 37913.
% \bibitem[Li et al.(2024)]{2024GCN.37844....1L} Li, W.~X. and 8 colleagues 2024.\ EP241021a: GSP detects optical candidate counterpart.\ GRB Coordinates Network, Circular Service, No. 37844 37844.
% \bibitem[Newsome et al.(2024)]{2024TNSCR3794....1N} Newsome, M., Andrews, M., \textbf{Farah, J.}, Howell, D.~A., McCully, C.\ 2024.\ Global SN Project Transient Classification Report for 2024-09-30.\ Transient Name Server Classification Report 2024-3794.
% \bibitem[Andrews et al.(2024)]{2024TNSCR3490....1A} Andrews, M., \textbf{Farah, J.}, Howell, D.~A., McCully, C.\ 2024.\ Transient Classification Report for 2024-09-13.\ Transient Name Server Classification Report 2024-3490.

% \bibitem[Andrews et al.(2024)]{2024TNSCR3471....1A} Andrews, M., \textbf{Farah, J.}, Howell, D.~A., McCully, C.\ 2024.\ Transient Classification Report for 2024-09-12.\ Transient Name Server Classification Report 2024-3471.
% \bibitem[Bostroem et al.(2024)]{2024TNSCR2571....1B} Bostroem, K.~A. and 24 colleagues 2024.\ DLT40 Transient Classification Report for 2024-07-24.\ Transient Name Server Classification Report 2024-2571.
% \bibitem[Newsome et al.(2024)]{2024TNSCR1687....1N} Newsome, M., \textbf{Farah, J.}, Howell, D.~A., McCully, C., Gonzalez, E.~P., Terreran, G.\ 2024.\ Global SN Project Transient Classification Report for 2024-05-28.\ Transient Name Server Classification Report 2024-1687.
% \bibitem[Hosseinzadeh et al.(2024)]{2024TNSCR1686....1H} Hosseinzadeh, G. and 8 colleagues 2024.\ Global SN Project Transient Classification Report for 2024-05-28.\ Transient Name Server Classification Report 2024-1686.
% \bibitem[Terreran et al.(2024)]{2024TNSCR.843....1T} Terreran, G., Howell, D.~A., McCully, C., Newsome, M., Gonzalez, E.~P., \textbf{Farah, J.}\ 2024.\ Global SN Project Transient Classification Report for 2024-03-28.\ Transient Name Server Classification Report 2024-843.
% \bibitem[Terreran et al.(2024)]{2024TNSCR.828....1T} Terreran, G., Howell, D.~A., McCully, C., Newsome, M., Gonzalez, E.~P., \textbf{Farah, J.}\ 2024.\ Global SN Project Transient Classification Report for 2024-03-27.\ Transient Name Server Classification Report 2024-828.
% \bibitem[Newsome et al.(2024)]{2024TNSCR.505....1N} Newsome, M. and 7 colleagues 2024.\ Global SN Project Transient Classification Report for 2024-02-21.\ Transient Name Server Classification Report 2024-505.
% \bibitem[Newsome et al.(2024)]{2024TNSCR.448....1N} Newsome, M., \textbf{Farah, J.}, Howell, D.~A., McCully, C., Gonzalez, E.~P., Terreran, G.\ 2024.\ Global SN Project Transient Classification Report for 2024-02-14.\ Transient Name Server Classification Report 2024-448.
% \bibitem[Newsome et al.(2024)]{2024TNSCR.413....1N} Newsome, M., \textbf{Farah, J.}, Howell, D.~A., McCully, C., Gonzalez, E.~P., Terreran, G.\ 2024.\ Global SN Project Transient Classification Report for 2024-02-12.\ Transient Name Server Classification Report 2024-413.
% \bibitem[Gonzalez et al.(2024)]{2024TNSCR.399....1G} Gonzalez, E.~P. and 6 colleagues 2024.\ Global SN Project Transient Classification Report for 2024-02-11.\ Transient Name Server Classification Report 2024-399.
% \bibitem[Newsome et al.(2024)]{2024TNSCR.344....1N} Newsome, M. and 8 colleagues 2024.\ Global SN Project Transient Classification Report for 2024-02-04.\ Transient Name Server Classification Report 2024-344.
% \bibitem[Newsome et al.(2024)]{2024TNSCR.261....1N} Newsome, M. and 6 colleagues 2024.\ Global SN Project Transient Classification Report for 2024-01-26.\ Transient Name Server Classification Report 2024-261.
% \bibitem[Terreran et al.(2024)]{2024TNSCR.232....1T} Terreran, G. and 6 colleagues 2024.\ Global SN Project Transient Classification Report for 2024-01-22.\ Transient Name Server Classification Report 2024-232.
% \bibitem[Newsome et al.(2024)]{2024TNSCR..44....1N} Newsome, M. and 6 colleagues 2024.\ Global SN Project Transient Classification Report for 2024-01-05.\ Transient Name Server Classification Report 2024-44.
% \bibitem[Arcavi et al.(2023)]{2023TNSCR3391....1A} Arcavi, I., Terreran, G., Newsome, M., \textbf{Farah, J.}, Charalampopoulos, P.\ 2023.\ StarDestroyers Transient Classification Report for 2023-12-30.\ Transient Name Server Classification Report 2023-3391.
% \bibitem[Terreran et al.(2023)]{2023TNSCR3376....1T} Terreran, G., Howell, D.~A., McCully, C., Newsome, M., Gonzalez, E.~P., \textbf{Farah, J.}\ 2023.\ Global SN Project Transient Classification Report for 2023-12-28.\ Transient Name Server Classification Report 2023-3376.
% \bibitem[Farah et al.(2023)]{2023TNSCR2992....1F} \textbf{Farah, J.} and 6 colleagues 2023.\ Transient Classification Report for 2023-11-17.\ Transient Name Server Classification Report 2023-2992.
% \bibitem[Newsome et al.(2023)]{2023TNSCR2557....1N} Newsome, M., \textbf{Farah, J.}, Howell, D.~A., McCully, C., Gonzalez, E.~P., Terreran, G.\ 2023.\ Global SN Project Transient Classification Report for 2023-10-11.\ Transient Name Server Classification Report 2023-2557.
% \bibitem[Gonzalez et al.(2023)]{2023TNSCR2348....1G} Gonzalez, E.~P. and 6 colleagues 2023.\ Global SN Project Transient Classification Report for 2023-09-21.\ Transient Name Server Classification Report 2023-2348.
% \bibitem[Pellegrino et al.(2023)]{2023TNSCR2244....1P} Pellegrino, C. and 6 colleagues 2023.\ Global SN Project Transient Classification Report for 2023-09-12.\ Transient Name Server Classification Report 2023-2244.
% \bibitem[Farah et al.(2023)]{2023TNSCR2213....1F} \textbf{Farah, J.} and 6 colleagues 2023.\ Transient Classification Report for 2023-09-08.\ Transient Name Server Classification Report 2023-2213.
% \bibitem[Terreran et al.(2023)]{2023TNSCR2212....1T} Terreran, G. and 6 colleagues 2023.\ Transient Classification Report for 2023-09-08.\ Transient Name Server Classification Report 2023-2212.
% \bibitem[Terreran et al.(2023)]{2023TNSCR2162....1T} Terreran, G. and 6 colleagues 2023.\ Global SN Project Transient Classification Report for 2023-09-04.\ Transient Name Server Classification Report 2023-2162.
% \bibitem[Pellegrino et al.(2023)]{2023TNSCR2134....1P} Pellegrino, C. and 6 colleagues 2023.\ Global SN Project Transient Classification Report for 2023-08-31.\ Transient Name Server Classification Report 2023-2134.
% \bibitem[Gonzalez et al.(2023)]{2023TNSCR2054....1G} Gonzalez, E.~P. and 6 colleagues 2023.\ Global SN Project Transient Classification Report for 2023-08-23.\ Transient Name Server Classification Report 2023-2054.
% \bibitem[Newsome et al.(2023)]{2023TNSCR2033....1N} Newsome, M. and 6 colleagues 2023.\ Global SN Project Transient Classification Report for 2023-08-21.\ Transient Name Server Classification Report 2023-2033.
% \bibitem[Newsome et al.(2023)]{2023TNSCR1999....1N} Newsome, M. and 6 colleagues 2023.\ Global SN Project Transient Classification Report for 2023-08-17.\ Transient Name Server Classification Report 2023-1999.
% \bibitem[Farah et al.(2023)]{2023TNSCR1960....1F} \textbf{Farah, J.} and 6 colleagues 2023.\ Transient Classification Report for 2023-08-14.\ Transient Name Server Classification Report 2023-1960.
% \bibitem[Newsome et al.(2023)]{2023TNSCR1865....1N} Newsome, M. and 8 colleagues 2023.\ Global SN Project Transient Classification Report for 2023-08-05.\ Transient Name Server Classification Report 2023-1865.
% \bibitem[Newsome et al.(2023)]{2023TNSCR1767....1N} Newsome, M. and 7 colleagues 2023.\ Global SN Project Transient Classification Report for 2023-07-25.\ Transient Name Server Classification Report 2023-1767.
% \bibitem[Newsome et al.(2023)]{2023TNSCR1754....1N} Newsome, M. and 8 colleagues 2023.\ Global SN Project Transient Classification Report for 2023-07-24.\ Transient Name Server Classification Report 2023-1754.
% \bibitem[Gonzalez et al.(2023)]{2023TNSCR1753....1G} Gonzalez, E.~P. and 6 colleagues 2023.\ Global SN Project Transient Classification Report for 2023-07-24.\ Transient Name Server Classification Report 2023-1753.
% \bibitem[Newsome et al.(2023)]{2023TNSCR1737....1N} Newsome, M. and 6 colleagues 2023.\ Global SN Project Transient Classification Report for 2023-07-22.\ Transient Name Server Classification Report 2023-1737.
% \bibitem[Gonzalez et al.(2023)]{2023TNSCR1496....1G} Gonzalez, E.~P. and 6 colleagues 2023.\ Global SN Project Transient Classification Report for 2023-06-26.\ Transient Name Server Classification Report 2023-1496.
% \bibitem[Pellegrino et al.(2023)]{2023TNSCR1474....1P} Pellegrino, C. and 6 colleagues 2023.\ Global SN Project Transient Classification Report for 2023-06-23.\ Transient Name Server Classification Report 2023-1474.
% \bibitem[Pellegrino et al.(2023)]{2023TNSCR1279....1P} Pellegrino, C. and 7 colleagues 2023.\ Global SN Project Transient Classification Report for 2023-05-30.\ Transient Name Server Classification Report 2023-1279.
% \bibitem[Farah et al.(2023)]{2023TNSCR1137....1F} \textbf{Farah, J.} and 6 colleagues 2023.\ Transient Classification Report for 2023-05-17.\ Transient Name Server Classification Report 2023-1137.
% \bibitem[Gonzalez et al.(2023)]{2023TNSCR1017....1G} Gonzalez, E.~P. and 6 colleagues 2023.\ Global SN Project Transient Classification Report for 2023-05-05.\ Transient Name Server Classification Report 2023-1017.
% \bibitem[Newsome et al.(2023)]{2023TNSCR.944....1N} Newsome, M. and 6 colleagues 2023.\ Global SN Project Transient Classification Report for 2023-04-27.\ Transient Name Server Classification Report 2023-944.
% \bibitem[Farah et al.(2023)]{2023TNSCR.793....1F} \textbf{Farah, J.} and 6 colleagues 2023.\ Transient Classification Report for 2023-04-14.\ Transient Name Server Classification Report 2023-793.
% \bibitem[Pellegrino et al.(2023)]{2023TNSCR.725....1P} Pellegrino, C. and 6 colleagues 2023.\ Global SN Project Transient Classification Report for 2023-04-06.\ Transient Name Server Classification Report 2023-725.
% \bibitem[Pellegrino et al.(2023)]{2023TNSAN..75....1P} Pellegrino, C. and 6 colleagues 2023.\ Classification of AT 2023emq as a Type Icn Supernova.\ Transient Name Server AstroNote 75.
% \bibitem[Pellegrino et al.(2023)]{2023TNSCR.542....1P} Pellegrino, C. and 7 colleagues 2023.\ Global SN Project Transient Classification Report for 2023-03-15.\ Transient Name Server Classification Report 2023-542.
% \bibitem[Terreran et al.(2023)]{2023TNSCR.490....1T} Terreran, G. and 6 colleagues 2023.\ Global SN Project Transient Classification Report for 2023-03-07.\ Transient Name Server Classification Report 2023-490.
% \bibitem[Hosseinzadeh et al.(2023)]{2023TNSCR.277....1H} Hosseinzadeh, G. and 8 colleagues 2023.\ Global SN Project Transient Classification Report for 2023-02-03.\ Transient Name Server Classification Report 2023-277.
% \bibitem[Pellegrino et al.(2023)]{2023TNSCR.218....1P} Pellegrino, C. and 6 colleagues 2023.\ Global SN Project Transient Classification Report for 2023-01-27.\ Transient Name Server Classification Report 2023-218.
% \bibitem[Farah et al.(2023)]{2023TNSCR.144....1F} \textbf{Farah, J.} and 6 colleagues 2023.\ Transient Classification Report for 2023-01-20.\ Transient Name Server Classification Report 2023-144.
% \bibitem[Pellegrino et al.(2022)]{2022TNSCR2757....1P} Pellegrino, C. and 8 colleagues 2022.\ Global SN Project Transient Classification Report for 2022-09-23.\ Transient Name Server Classification Report 2022-2757.
% \bibitem[Pellegrino et al.(2022)]{2022TNSCR2707....1P} Pellegrino, C. and 7 colleagues 2022.\ Global SN Project Transient Classification Report for 2022-09-19.\ Transient Name Server Classification Report 2022-2707.
% \bibitem[Pellegrino et al.(2022)]{2022TNSCR2535....1P} Pellegrino, C. and 8 colleagues 2022.\ Global SN Project Transient Classification Report for 2022-08-31.\ Transient Name Server Classification Report 2022-2535.
% \bibitem[Pellegrino et al.(2022)]{2022TNSCR2510....1P} Pellegrino, C. and 7 colleagues 2022.\ Global SN Project Transient Classification Report for 2022-08-29.\ Transient Name Server Classification Report 2022-2510.

% \section{\sc Journal Publications}
% \vspace{+.4cm}
% \section{\sc Invited Talks}
% \vspace{+.4cm}
% \section{\sc Teaching}
% \vspace{+.4cm}

\end{resume}

\end{document}